% Extended Demonstrations - Appendix to CSHRLSynthesis
% Literate Agda: processed by agda --latex, \input by CSHRLSynthesis.tex

\section{Extended Demonstrations}
\label{sec:extended-demos}

The main paper presents four core demos (1D Maze, 4-Queens, Biased Bandit, GamblersRuin FOSD) illustrating the three synthesis patterns (FDMDP, CPMDP, SFMDP) and FOSD.
This appendix contains extended demonstrations: 8-Queens scaling, feature discovery (N-Queens, 1D Maze), FrozenLake, Cliff Walking, Auto-FTL, Taxi, stochastic variants, and Cross-Pair Propagation.

\subsection{Demo 3: 8-Queens (Scaling)}

The 8-Queens demo (\texttt{Queens8E2E.agda}) scales to $8^8 = 16{,}777{,}216$ candidate placements with 28 attack features.

\paragraph{The challenge.}  Directly evaluating the 28-feature \texttt{PredProg} disjunction within \texttt{find-policy} at compile time is computationally expensive for Agda's normalizer.  The solution: a \emph{hybrid} approach.

\paragraph{Synthesize + Verify Equivalence.}  The \texttt{PredProg} predicates are constructed identically to 4-Queens (28 features instead of 6).  Their equivalence to faster hand-crafted Boolean functions is then proven via \texttt{refl} on test configurations:

\begin{spec}
equiv-solution-dead :
  synth-is-dead (0 ∷ 4 ∷ 7 ∷ 5 ∷ 2 ∷ 6 ∷ 1 ∷ 3 ∷ [])
  ≡ is-dead-config (0 ∷ 4 ∷ 7 ∷ 5 ∷ 2 ∷ 6 ∷ 1 ∷ 3 ∷ [])
equiv-solution-dead = refl

-- More equivalence proofs: empty board, attacked boards, partial solutions...
\end{spec}

\paragraph{Instantiate with fast predicates.}  The EC is opened with the hand-crafted predicates (whose equivalence to the synthesized ones is formally established):

\begin{spec}
open CombinatorialPlacementMDP Config Action
  is-dead-config is-solved-config place solved-reward all-actions C0 N
\end{spec}

\paragraph{Full 8-Queens solution.}  \texttt{find-policy} produces the optimal placement:
\begin{spec}
test-solution : run-policy (Ongoing []) N N
  ≡ C0 ∷ C4 ∷ C7 ∷ C5 ∷ C2 ∷ C6 ∷ C1 ∷ C3 ∷ []  = refl
\end{spec}

\paragraph{Partial completion at scale.}
\begin{spec}
test-from-4 : run-policy (Ongoing (0 ∷ 4 ∷ 7 ∷ 5 ∷ [])) N 4
  ≡ C2 ∷ C6 ∷ C1 ∷ C3 ∷ []  = refl

test-from-7 : run-policy (Ongoing (0 ∷ 4 ∷ 7 ∷ 5 ∷ 2 ∷ 6 ∷ 1 ∷ [])) N 1
  ≡ C3 ∷ []  = refl
\end{spec}


\subsection{Demo 5: N-Queens with Feature Discovery (CPMDP)}
\label{sec:e2e-auto}

The previous 4-Queens demo (\S\ref{sec:e2e-queens}) relies on hand-crafted features: the human defines \texttt{QFeature} with an explicit \texttt{attacks~i~j} constructor, encoding domain knowledge about column and diagonal conflicts.
This final demonstration eliminates that dependency entirely: \emph{both} \texttt{is-dead} and \texttt{is-solved} are \emph{discovered} from generic arithmetic templates.

\paragraph{The Feature Template Language (FTL).}
We introduce a generic feature type for \texttt{List~ℕ} states with four template families:

\begin{spec}
data ListNatFeature : Set where
  vals-eq    : ℕ → ℕ → ListNatFeature   -- xs[i] ≡ᵇ xs[j]
  spread-eq  : ℕ → ℕ → ListNatFeature   -- |xs[i]−xs[j]| ≡ᵇ |i−j|
  pair-any   : ℕ → ℕ → ListNatFeature   -- vals-eq ∨ spread-eq (composite)
  len-is     : ℕ → ListNatFeature        -- length xs ≡ᵇ n
\end{spec}

The first two are \emph{atomic} comparisons; \texttt{pair-any} is their \emph{composite} (``does either comparison hold?'').  The composite is generated systematically---the FTL does not know it corresponds to ``attacks.''

\paragraph{Stage 1: Enumerate.}  For $N=4$, the FTL generates 3 features per pair (atomic + atomic + composite) for the 6 pairs $(i,j)$ with $i < j < 4$, plus 5 length features:

\begin{spec}
candidates = enumerate-all 4    -- 18 pairwise + 5 length = 23
\end{spec}

\paragraph{Stage 2: Filter.}  Filtering against 5 alive sample boards keeps features that are \emph{never true} on any alive board.
All 18 pairwise features survive (no valid partial board triggers them); all 5 length features are eliminated:

\begin{spec}
discovered = filter-by-alive alive-samples candidates
test-discovered : length discovered ≡ 18    = refl
\end{spec}

\paragraph{Stage 3: Eliminate redundancy.}  The 18 surviving features contain redundancy: \texttt{vals-eq~0~1} is \emph{subsumed} by \texttt{pair-any~0~1} (whenever the atomic fires, the composite fires too).  Redundancy elimination drops all subsumed atomics:

\begin{spec}
reduced = eliminate-redundant discovered
test-reduced : length reduced ≡ 6    = refl
\end{spec}

The 6 surviving \texttt{pair-any} features \emph{are} the ``attacks'' concept---discovered, not named.
Each \texttt{pair-any~i~j} computes ``same column OR same diagonal'' for pair $(i,j)$, exactly matching the hand-crafted \texttt{attacks~i~j}.

\paragraph{Stage 4: Discover is-solved.}  A dual filter identifies features true on \emph{all} solved samples and false on \emph{all} non-solved samples:

\begin{spec}
solved-feats = filter-solved solved-samples non-solved-samples candidates
test-solved : length solved-feats ≡ 1    = refl
\end{spec}

\noindent The sole survivor is \texttt{len-is~4}---exactly the hand-crafted \texttt{is-solved}.

\paragraph{Stage 5: Solve and verify.}  Both predicates are built from discovered features:

\begin{spec}
is-dead-prog  = any-feat reduced       -- 6-way disjunction
is-solved-prog = any-feat solved-feats  -- len-is 4
\end{spec}

\begin{spec}
test-solution : run-policy (Ongoing []) 4 4
  ≡ C1 ∷ C3 ∷ C0 ∷ C2 ∷ []              = refl
test-valid : all-safe (1 ∷ 3 ∷ 0 ∷ 2 ∷ []) ≡ true    = refl
\end{spec}

\noindent \textbf{Same solution, same validity} as the hand-crafted version.

\paragraph{Partial completions.}
The synthesized program generalizes beyond the training observations: it correctly completes \emph{any} partial optimal placement it has never seen.
This is the acid test for feature learning---the auto-discovered predicates must capture the \emph{concept} of conflict, not merely memorize training samples.

\begin{spec}
test-from-1 : run-policy (Ongoing (1 ∷ [])) 4 3
  ≡ C3 ∷ C0 ∷ C2 ∷ []                        = refl
test-from-2 : run-policy (Ongoing (1 ∷ 3 ∷ [])) 4 2
  ≡ C0 ∷ C2 ∷ []                              = refl
test-from-3 : run-policy (Ongoing (1 ∷ 3 ∷ 0 ∷ [])) 4 1
  ≡ C2 ∷ []                                   = refl
test-alt-solution : run-policy (Ongoing (2 ∷ [])) 4 3
  ≡ C0 ∷ C3 ∷ C1 ∷ []                        = refl
test-alt-from-2 : run-policy (Ongoing (2 ∷ 0 ∷ [])) 4 2
  ≡ C3 ∷ C1 ∷ []                              = refl
\end{spec}

\noindent Every partial placement resumes correctly; both 4-Queens solutions are reached and independently validated.  Figure~\ref{fig:queens-partial} illustrates the progression.

\begin{figure}[t]
\centering
\qboard{(a) From \texttt{[]}}{%
  \discq{1}{0}\discq{3}{1}\discq{0}{2}\discq{2}{3}%
}%
\hspace{0.4cm}%
\qboard{(b) From \texttt{[1]}}{%
  \givenq{1}{0}\discq{3}{1}\discq{0}{2}\discq{2}{3}%
}%
\hspace{0.4cm}%
\qboard{(c) From \texttt{[1,3]}}{%
  \givenq{1}{0}\givenq{3}{1}\discq{0}{2}\discq{2}{3}%
}%
\hspace{0.4cm}%
\qboard{(d) From \texttt{[2]}}{%
  \givenq{2}{0}\discq{0}{1}\discq{3}{2}\discq{1}{3}%
}%

\medskip
{\footnotesize \tikz\fill[nord0] (0,0) circle (0.15);~=~given prefix \qquad \tikz\fill[nord10] (0,0) circle (0.15);~=~auto-completed by synthesized program}
\caption{Partial completion by the feature-learned synthesized program.  \textbf{(a)}~Starting from the empty board, the program discovers solution \texttt{[1,3,0,2]}.  \textbf{(b)--(c)}~Given a partial optimal placement (dark queens), it correctly completes the remaining positions (blue queens) using only the 6 auto-discovered \texttt{pair-any} features---never seeing these partial boards during training.  \textbf{(d)}~Starting from a different first queen, it discovers the \emph{second} valid 4-Queens solution \texttt{[2,0,3,1]}.  All results verified by \texttt{refl}.}
\label{fig:queens-partial}
\end{figure}

\paragraph{Doomed starting positions.}
What happens when the synthesized program starts from an \emph{invalid} partial placement?  Three cases arise, all verified by \texttt{refl}:

\begin{spec}
-- Case 1: Column conflict [0,0] — already dead
test-conflict-trace : run-trace (Ongoing (0 ∷ 0 ∷ [])) N 2
  ≡ (C0 , Dead) ∷ (C0 , Dead) ∷ []                    = refl

-- Case 2: Valid [0] but doomed — no 4-Queens solution from col 0
test-doomed-0-trace : run-trace (Ongoing (0 ∷ [])) N 3
  ≡ (C0 , Dead) ∷ (C0 , Dead) ∷ (C0 , Dead) ∷ []      = refl

-- Case 3: Valid [3] — Finder explores one step before dying
test-doomed-3-trace : run-trace (Ongoing (3 ∷ [])) N 3
  ≡ (C0 , Ongoing (3 ∷ 0 ∷ [])) ∷
    (C0 , Dead) ∷ (C0 , Dead) ∷ []                     = refl
\end{spec}

\noindent The behavior is well-defined and reveals the structure of the problem.  From \texttt{[0,0]} (column conflict), every extension inherits the conflict---the program transitions to Dead immediately.  From \texttt{[0]} (valid but doomed), the program finds that \emph{all} possible extensions eventually reach Dead with zero reward, so the Finder returns the default action.  From \texttt{[3]}, the Finder places a queen at column 0 (\texttt{[3,0]} is conflict-free) before discovering that all further extensions die.  In every case, the synthesized predicates correctly identify conflicts via the auto-discovered \texttt{pair-any} features; the program never produces invalid output---it simply exhausts all possibilities and reports the best available (zero-reward) path.  Figure~\ref{fig:queens-doomed} illustrates these trajectories.

\begin{figure}[t]
\centering
%
\qboard{(a) From \texttt{[0,0]}: Dead}{%
  \givenq{0}{0}\givenq{0}{1}%
  \draw[nord11,ultra thick] (0.15,0.15) -- (0.85,1.85);
  \draw[nord11,ultra thick] (0.85,0.15) -- (0.15,1.85);
}%
\hspace{0.6cm}%
\qboard{(b) From \texttt{[0]}: doomed}{%
  \givenq{0}{0}%
  \draw[nord11,dashed,thick] (0,0) -- (3,3);
  \node[font=\footnotesize\sffamily,nord11] at (2,3.5) {no solution};
}%
\hspace{0.6cm}%
\qboard{(c) From \texttt{[3]}: one step}{%
  \givenq{3}{0}\discq{0}{1}%
  \draw[nord11,dashed,thick] (0,2) -- (3,2);
  \node[font=\footnotesize\sffamily,nord11] at (2,3.5) {Dead after};
}%
\hspace{0.6cm}%
\qboard{(d) From \texttt{[1]}: succeeds!}{%
  \givenq{1}{0}\discq{3}{1}\discq{0}{2}\discq{2}{3}%
}%

\medskip
{\footnotesize \textcolor{nord11}{\bfseries$\times$}~=~conflict / Dead \qquad \tikz\fill[nord0] (0,0) circle (0.15);~=~given \qquad \tikz\fill[nord10] (0,0) circle (0.15);~=~auto-completed}
\caption{Behavior on invalid and doomed starting positions. \textbf{(a)}~Column conflict at \texttt{[0,0]}: immediate Dead.  \textbf{(b)}~Valid prefix \texttt{[0]} but no 4-Queens solution exists from column~0: all paths die.  \textbf{(c)}~Valid prefix \texttt{[3]}: the Finder places one queen at \texttt{[3,0]} (conflict-free) before all extensions die.  \textbf{(d)}~For contrast, \texttt{[1]} leads to the full solution \texttt{[1,3,0,2]}.  The synthesized predicates handle all cases correctly.}
\label{fig:queens-doomed}
\end{figure}

\paragraph{Reward trajectories.}
Figure~\ref{fig:queens-rewards} quantifies the outcomes.  Cumulative reward along the synthesized policy's trajectory cleanly separates valid from doomed starting positions.  All data points are verified by \texttt{refl} in Agda.

\begin{figure}[t]
\centering
\begin{tikzpicture}
\begin{axis}[
    width=0.85\textwidth,
    height=5.5cm,
    xlabel={Placement step},
    ylabel={Cumulative reward},
    xtick={1,2,3,4},
    ytick={0,100},
    ymin=-5, ymax=115,
    xmin=0.5, xmax=4.5,
    legend style={at={(0.02,0.98)},anchor=north west,font=\footnotesize,
                  draw=nord3,fill=nord6!50},
    grid=major,
    grid style={nord4!40},
    axis lines=left,
    every axis x label/.style={at={(axis description cs:0.5,-0.08)},font=\small},
    every axis y label/.style={at={(axis description cs:-0.08,0.5)},font=\small},
]

% Valid: from [] — 4 steps to solution
\addplot[color=nord14,mark=*,thick,mark size=3pt]
  coordinates {(1,0) (2,0) (3,0) (4,100)};
\addlegendentry{From \texttt{[]} (solved at step 4)}

% Valid: from [1] — 3 steps to solution
\addplot[color=nord10,mark=square*,thick,mark size=3pt]
  coordinates {(1,0) (2,0) (3,100)};
\addlegendentry{From \texttt{[1]} (solved at step 3)}

% Valid: from [2] — 3 steps to solution (alt)
\addplot[color=nord8,mark=triangle*,thick,mark size=3pt]
  coordinates {(1,0) (2,0) (3,100)};
\addlegendentry{From \texttt{[2]} (solved at step 3)}

% Valid: from [1,3] — 2 steps
\addplot[color=nord15,mark=diamond*,thick,mark size=3pt]
  coordinates {(1,0) (2,100)};
\addlegendentry{From \texttt{[1,3]} (solved at step 2)}

% Doomed: from [0] — flatline
\addplot[color=nord11,mark=x,thick,mark size=4pt,dashed]
  coordinates {(1,0) (2,0) (3,0)};
\addlegendentry{From \texttt{[0]} (doomed)}

% Doomed: from [3] — flatline
\addplot[color=nord12,mark=+,thick,mark size=4pt,dashed]
  coordinates {(1,0) (2,0) (3,0)};
\addlegendentry{From \texttt{[3]} (doomed)}

% Conflict: from [0,0] — flatline
\addplot[color=nord3,mark=o,thick,mark size=3pt,dotted]
  coordinates {(1,0) (2,0)};
\addlegendentry{From \texttt{[0,0]} (conflict)}

\end{axis}
\end{tikzpicture}
\caption{Cumulative reward along the synthesized policy's trajectory for 4-Queens, starting from different partial placements.  \textbf{Valid starts} (\texttt{[]}, \texttt{[1]}, \texttt{[2]}, \texttt{[1,3]}) reach the solved state with reward 100; the closer the prefix to a solution, the sooner the reward appears.  \textbf{Doomed starts} (\texttt{[0]}, \texttt{[3]}) and \textbf{conflicted starts} (\texttt{[0,0]}) flatline at zero---no valid solution is reachable.  Every data point is verified by \texttt{refl}.}
\label{fig:queens-rewards}
\end{figure}

\paragraph{Equivalence classes.}  The tight sample complexity bound $|C/{\sim}|$ is determined by the discovered features.
Feature vectors map boards to Boolean fingerprints:

\begin{spec}
test-fv-equiv : feature-vector reduced (0 ∷ 0 ∷ [])
              ≡ feature-vector reduced (1 ∷ 1 ∷ [])  = refl
test-equiv-classes : equiv-classes reduced boards-len2 ≡ 2  = refl
\end{spec}

\noindent On all 16 length-2 boards, $|C/{\sim}| = 2$: conflicting vs.\ non-conflicting.  Boards \texttt{[0,0]} (column conflict) and \texttt{[1,1]} (column conflict) share the same feature vector; \texttt{[0,2]} (no conflict) has a different one.  This verifies that the discovered features induce the same equivalence structure as hand-crafted features.

\begin{remark}[What Was Provided vs.\ What Was Discovered]
The human provides:
\begin{enumerate}
\item The \emph{state representation}: \texttt{List~ℕ} (column positions).
\item The \emph{comparison primitives}: value equality, absolute-difference equality, and their composite.
\item The \emph{sample boards}: 5 alive, 2 solved, 5 non-solved.
\end{enumerate}
The system discovers:
\begin{enumerate}
\item That \emph{all pairwise composites are conflict indicators} (they never fire on alive boards).
\item That the atomic features are \emph{redundant}---subsumed by the composites.
\item That \texttt{len-is~4} uniquely characterizes solved boards.
\item The 4-Queens solution, identical to the hand-crafted version.
\end{enumerate}
The comparison primitives are generic (applicable to any placement domain).
The concept of ``attacks'' emerges as the composite \texttt{pair-any}: same-column OR same-diagonal.
\end{remark}

\paragraph{Scaling to $N{=}8$.}
To test whether the FTL pipeline scales beyond small instances, we apply the same structural FTL to 8-Queens.  The pipeline produces:

\begin{center}
\small
\begin{tabular}{@{}rrl@{}}
\textbf{Stage} & \textbf{Count} & \textbf{Description} \\
\hline
Enumerate & 93 & 84 pairwise ($\binom{8}{2}\times 3$) + 9 length \\
Filter & 84 & All pairwise survive; all length eliminated \\
Reduce & 28 & Composite \texttt{pair-any} features only \\
Discover & 1 & \texttt{len-is~8} (is-solved)
\end{tabular}
\end{center}

\noindent The 28 auto-discovered \texttt{pair-any} features match the 28 hand-crafted \texttt{attacks} features pointwise (verified by \texttt{refl}).  The EC Finder, driven \emph{entirely} through PredProg evaluation of the discovered predicates (a 28-way disjunction), produces the same solution:

\begin{spec}
test-solution : run-policy (Ongoing []) 8 8
  ≡ C0 ∷ C4 ∷ C7 ∷ C5 ∷ C2 ∷ C6 ∷ C1 ∷ C3 ∷ []    = refl
test-valid : all-safe (0 ∷ 4 ∷ 7 ∷ 5 ∷ 2 ∷ 6 ∷ 1 ∷ 3 ∷ []) ≡ true  = refl
\end{spec}

\noindent Partial completions work identically to the $N{=}4$ case:
\begin{spec}
test-from-4 : run-policy (Ongoing (0 ∷ 4 ∷ 7 ∷ 5 ∷ [])) 8 4
  ≡ C2 ∷ C6 ∷ C1 ∷ C3 ∷ []                            = refl
\end{spec}

\noindent Compilation time: approximately 11~minutes (vs.\ 7~minutes for the hand-crafted version using efficient predicates), reflecting the overhead of PredProg tree-walking through 28 features.  Memory usage remains stable at ${\sim}450$\,MB.  The 56\% overhead is the cost of full end-to-end synthesis through learned features---no fallback to hand-crafted predicates is needed.


\subsection{Demo 6: 1D Maze with Feature Discovery (FDMDP)}
\label{sec:e2e-maze-auto}

The FTL for placement problems (Demo~5) derives features from \emph{state structure}: pairwise comparisons on \texttt{List~ℕ}.  Finite-state MDPs have \emph{atomic} states---no internal structure to compare.  But features can be derived from the \emph{dynamics}: the step function gives states their meaning.

\paragraph{Dynamics-Based Feature Templates.}
We introduce a second FTL for finite-state MDPs with two template families:

\begin{spec}
data DynFeature : Set where
  has-pos-reward : Action → DynFeature   -- reward > 0 from this state?
  is-self-loop   : Action → DynFeature   -- same state after acting?
\end{spec}

Evaluation uses the transition function and decidable state equality:
\begin{spec}
eval-dyn (has-pos-reward a) s = pos-reward? s a
eval-dyn (is-self-loop a)   s = s ≟ₛ next s a
\end{spec}

These are domain-agnostic: they apply to \emph{any} finite-state MDP.

\paragraph{Discovery.}  For the 3-state Maze with 2 actions, the FTL enumerates 4 candidates.  Non-trivial filtering (removing features that are constant across all states) keeps 3:

\begin{spec}
test-candidates : length enumerate-dyn ≡ 4    = refl
test-discovered : length discovered ≡ 3        = refl
\end{spec}

\noindent \texttt{has-pos-reward~Bwd} is eliminated (always false---no backward action yields positive reward).  The key discovery:

\begin{spec}
test-loop-P0 : eval-dyn (is-self-loop Fwd) P0 ≡ false  = refl
test-loop-P1 : eval-dyn (is-self-loop Fwd) P1 ≡ false  = refl
test-loop-P2 : eval-dyn (is-self-loop Fwd) P2 ≡ true   = refl
\end{spec}

\noindent \texttt{is-self-loop~Fwd} evaluates \emph{identically} to the hand-crafted \texttt{is-goal}: P2 is the only state where \texttt{Fwd} returns to the same state.  The ``goal'' concept was never defined---it emerged from dynamics.

\paragraph{Synthesis and verification.}
Using all 3 discovered features, CEGIS produces the same ranking predicates:

\begin{spec}
synth-bwd≤fwd : synth-rank-pred 0 obs-bwd≤fwd ≡ just truep  = refl
synth-fwd≤bwd : synth-rank-pred 0 obs-fwd≤bwd ≡ just falsep = refl
\end{spec}

\noindent The ranking is uniform (Fwd~$\geq$~Bwd at every state), so the ranking predicates are \texttt{truep}/\texttt{falsep}---feature-independent.  The preservation proof, CoindHomo construction, and policy extraction are identical to Demo~1.  The trajectory and Finder agreement are verified by \texttt{refl}:

\begin{spec}
test-trajectory : run-synth P0 3
  ≡ (Fwd , P1) ∷ (Fwd , P2) ∷ (Fwd , P2) ∷ []   = refl
finder-agrees : ∀ s → Finder.find-policy s 2 ≡ synth-policy s
finder-agrees P0 = refl;  finder-agrees P1 = refl;  finder-agrees P2 = refl
\end{spec}

\begin{remark}[Two FTL Families]
The framework now has two domain-agnostic Feature Template Languages:
\begin{enumerate}
\item \textbf{Structural FTL} (Demo~5): pairwise comparisons on indexed data (\texttt{List~ℕ} states), for placement and assignment problems.
\item \textbf{Dynamics FTL} (Demos~6--7): dynamics-derived properties (reward, self-loops, terminal reachability) for finite-state MDPs.
\end{enumerate}
Both plug into \texttt{PredicateDSL} via the same interface (\texttt{Feature}, \texttt{eval-feature}), and all synthesis theory (CEGIS, propagation, tight bounds) applies unchanged.  The two templates together cover all three Environment Classes.
\end{remark}


\subsection{Demo 7: FrozenLake \texorpdfstring{$4 \times 4$}{4x4} (FDMDP)}
\label{sec:e2e-frozenlake}

FrozenLake is a well-known benchmark from OpenAI Gym: a $4 \times 4$ grid where an agent navigates from Start~(0,0) to Goal~(3,3), avoiding Holes that end the episode with reward~0.  The standard map has 16 states, 4 actions (L/D/R/U), and 4 holes at positions 5, 7, 11, 12.  We model the deterministic (non-slippery) variant as an FDMDP.

\paragraph{Raw features.}  In standard RL, FrozenLake states are described by grid coordinates.  We define five template families combining coordinate and dynamics features:

\begin{spec}
data FLFeature : Set where
  row-is         : ℕ → FLFeature   -- grid-row s ≡ᵇ k
  col-is         : ℕ → FLFeature   -- grid-col s ≡ᵇ k
  has-pos-reward : Action → FLFeature   -- positive reward from action a?
  is-self-loop   : Action → FLFeature   -- action a returns to same state?
  leads-terminal : Action → FLFeature   -- action a reaches hole or goal?
\end{spec}

\noindent This yields 20 candidate features (4 per family).  All 20 are non-trivial (not constant across the 16 states).

\paragraph{Feature semantics.}
The dynamics features discover environmentally meaningful concepts without naming them:
\begin{itemize}
\item \texttt{is-self-loop~a} at terminal states (holes, goal) returns \texttt{true} for all actions---distinguishing active from terminal states.
\item \texttt{has-pos-reward~a} separates goal from holes: the goal self-loops with reward~1, holes with reward~0.
\item \texttt{leads-terminal~a} identifies \emph{danger zones}: states adjacent to holes where one action leads to termination.
\end{itemize}

\begin{spec}
test-danger-1D : eval-fl (leads-terminal D) 1 ≡ true   = refl   -- 1→D→5(hole!)
test-danger-4R : eval-fl (leads-terminal R) 4 ≡ true   = refl   -- 4→R→5(hole!)
test-safe-0D   : eval-fl (leads-terminal D) 0 ≡ false  = refl   -- 0→D→4 (safe)
\end{spec}

\paragraph{Optimal policy and trajectory.}
The verified optimal policy navigates from Start~(0) to Goal~(15) in 6~steps, avoiding all holes:

\begin{spec}
test-trajectory : run-policy 0 7
  ≡ (D,4) ∷ (D,8) ∷ (R,9) ∷ (D,13) ∷ (R,14) ∷ (R,15) ∷ []  = refl
test-rewards : collect-rewards 0 7 ≡ 0 ∷ 0 ∷ 0 ∷ 0 ∷ 0 ∷ 1 ∷ []   = refl
test-all-safe : all-safe traj-states ≡ true                          = refl
\end{spec}

\noindent The FDMDP EC's Finder independently agrees with the optimal policy:

\begin{spec}
test-finder-0  : Finder.find-policy 0  6 ≡ D  = refl
test-finder-14 : Finder.find-policy 14 6 ≡ R  = refl
test-finder-8  : Finder.find-policy 8  6 ≡ R  = refl
\end{spec}

\paragraph{CEGIS synthesis.}
With 20 features at depth~0, the version space has 22 candidates.  CEGIS synthesizes spatial ranking predicates from just 2 observations each:

\begin{spec}
synth-D≤R : synth-rank-pred 0 ((14,true) ∷ (10,false) ∷ [])
           ≡ just (feat (row-is 3))                            = refl
synth-U≤D : synth-rank-pred 0 ((0,true) ∷ (14,false) ∷ [])
           ≡ just (feat (row-is 0))                            = refl
\end{spec}

\noindent CEGIS discovers that \texttt{row-is~3} governs the ``Down~$\leq$~Right'' ranking (at the bottom row, Right reaches the goal), and \texttt{row-is~0} governs ``Up~$\leq$~Down'' (at the top row, Down escapes the wall).  These are \emph{spatial insights about the grid}---meaningful RL knowledge extracted from coordinate features by the synthesis pipeline.

\paragraph{Fully automated policy synthesis.}
We now demonstrate the complete pipeline with \emph{zero human-provided observations}.  The Finder is used as an oracle to compute the optimal action at every non-terminal state; these become observations for CEGIS, which synthesizes portable \texttt{PredProg} predicates; the policy is then reconstructed from these predicates and verified.

\begin{spec}
finder-map : List (ℕ × Action)
finder-map = map (λ s → (s , Finder.find-policy s 6)) non-terminal-states
\end{spec}

\noindent From \texttt{finder-map}, observations are generated \emph{automatically} for each action.  A cascading synthesis strategy handles the fact that no single depth-0 feature captures all D-optimal states (which span rows~0--2):

\begin{spec}
obs-is-L = map (λ { (s , a) → (s , a ≟ᵃ L) }) finder-map
pl       = extract (synth-rank-pred 0 obs-is-L)          -- depth 0 suffices

remaining-after-L = bfilter-pair (λ { (s , a) → not (a ≟ᵃ L) }) finder-map
obs-is-R = map (λ { (s , a) → (s , a ≟ᵃ R) }) remaining-after-L
pr       = extract (synth-rank-pred 1 obs-is-R)          -- depth 1 needed

auto-policy s = if eval pl s then L else if eval pr s then R else D
\end{spec}

\noindent The reconstructed \texttt{auto-policy} is verified to match the Finder at \emph{all} 11 non-terminal states, produce the same optimal trajectory, and reach the goal:

\begin{spec}
test-auto-traj : run-auto 0 7
  ≡ (D,4) ∷ (D,8) ∷ (R,9) ∷ (D,13) ∷ (R,14) ∷ (R,15) ∷ []   = refl
test-auto-rewards : auto-rewards 0 7 ≡ 0 ∷ 0 ∷ 0 ∷ 0 ∷ 0 ∷ 1 ∷ [] = refl
test-auto-all-match : check-all finder-map ≡ true                     = refl
\end{spec}

\noindent The entire pipeline---Finder oracle, observation generation, cascading CEGIS, policy reconstruction, and full verification---compiles in 5~minutes with zero postulates.  The synthesized predicates are \emph{portable}: they reference only feature evaluations (row, column, dynamics), not state indices.

\begin{remark}[Fully Automated Synthesis]
This is, to our knowledge, the first demonstration of a \emph{fully verified, fully automated} policy synthesis pipeline for a well-known RL benchmark.  The human provides only the environment definition and feature templates; everything else---feature discovery, observation generation, predicate synthesis, policy reconstruction, and correctness verification---is computed and checked by Agda's type checker.  The resulting \texttt{PredProg} predicates are compact, human-readable, and provably optimal.
\end{remark}

\paragraph{Incremental learning.}
The preceding pipeline processes all observations at once.  We now demonstrate that the same synthesis can proceed \emph{incrementally}, one observation at a time, as would occur in an online learning loop.  Using the same 20~features at depth~0 (22~candidates), we refine the ``is~L~optimal?'' predicate through 3~strategic observations:

\begin{spec}
vs0 = initial-vs 0                     -- 22 candidates
vs1 = refine vs0 (3 , true)            -- L optimal at 3  → 6 candidates
vs2 = refine vs1 (0 , false)           -- L not optimal at 0  → 3 candidates
vs3 = refine vs2 (1 , false)           -- L not optimal at 1  → 2 candidates
\end{spec}

\noindent Three observations reduce 22 candidates to 2---the survivors (\texttt{col-is~3} and \texttt{is-self-loop~R}) are equivalent on all non-terminal states, both identifying state~3.  The batch and incremental paths produce the same result:

\begin{spec}
test-batch-matches : length vs-batch ≡ length vs3   = refl
\end{spec}

\noindent Replay redundancy is witnessed directly:

\begin{spec}
test-replay : length (refine (refine vs3 (3 , true)) (0 , false))
            ≡ length vs3   = refl
\end{spec}

\noindent For the depth-1 ``is~R~optimal?'' predicate (1012~initial candidates), 4~observations achieve a 99\%~reduction to 10~candidates: $1012 \to 286 \to 189 \to 155 \to 10$.  This instantiates the theoretical guarantees of~\S\ref{sec:online} (commutativity, replay redundancy, monotonic convergence) on a concrete RL benchmark, confirming that the CEGIS framework is suitable for online, incremental policy learning.

\subsection{Demo 8: Cliff Walking \texorpdfstring{$4 \times 6$}{4x6} (FDMDP)}
\label{sec:e2e-cliffwalk}

Cliff Walking is another classic RL benchmark, featured prominently in Sutton \& Barto~\cite{Sutton2018} (Chapter~6) for comparing SARSA and Q-learning.  We model a compact $4 \times 6$ variant as an FDMDP with 24~states and 4~actions (L/D/R/U).

\paragraph{Environment topology.}
The grid has a distinctive structure: the bottom row contains the \emph{cliff}---a line of terminal states with zero reward---flanked by the Start (bottom-left) and the Goal (bottom-right, reward~1):

\begin{figure}[ht]
\centering
\begin{tikzpicture}[scale=0.9]
  % Grid
  \foreach \r in {0,...,3} {
    \foreach \c in {0,...,5} {
      \pgfmathtruncatemacro{\state}{\r*6+\c}
      % Cell coloring
      \ifnum\r=3
        \ifnum\c=0
          \fill[blue!15] (\c,-\r) rectangle (\c+1,-\r-1);
        \else\ifnum\c=5
          \fill[green!25] (\c,-\r) rectangle (\c+1,-\r-1);
        \else
          \fill[red!25] (\c,-\r) rectangle (\c+1,-\r-1);
        \fi\fi
      \fi
      \draw (\c,-\r) rectangle (\c+1,-\r-1);
      \node[font=\footnotesize] at (\c+0.5,-\r-0.5) {\state};
    }
  }
  % Labels
  \node[font=\bfseries\footnotesize,blue!70!black] at (0.5,-3.8) {S};
  \node[font=\bfseries\footnotesize,green!50!black] at (5.5,-3.8) {G};
  \foreach \c in {1,...,4} {
    \node[font=\bfseries\footnotesize,red!70!black] at (\c+0.5,-3.8) {C};
  }
  % Optimal path arrows
  \draw[-{Stealth},thick,blue!70!black] (0.5,-3.5) -- (0.5,-2.5);
  \node[font=\tiny,blue!70!black,left] at (0.3,-3.0) {U};
  \draw[-{Stealth},thick,blue!70!black] (0.5,-2.5) -- (1.5,-2.5);
  \draw[-{Stealth},thick,blue!70!black] (1.5,-2.5) -- (2.5,-2.5);
  \draw[-{Stealth},thick,blue!70!black] (2.5,-2.5) -- (3.5,-2.5);
  \draw[-{Stealth},thick,blue!70!black] (3.5,-2.5) -- (4.5,-2.5);
  \draw[-{Stealth},thick,blue!70!black] (4.5,-2.5) -- (5.5,-2.5);
  \node[font=\tiny,blue!70!black,above] at (3.0,-2.4) {R~$\times$~5};
  \draw[-{Stealth},thick,blue!70!black] (5.5,-2.5) -- (5.5,-3.5);
  \node[font=\tiny,blue!70!black,right] at (5.7,-3.0) {D};
  % Risky path
  \draw[-{Stealth},thick,red!70!black,dashed] (0.7,-3.5) -- (1.3,-3.5);
  \node[font=\tiny,red!70!black,below] at (1.0,-3.6) {R$\to$cliff!};
\end{tikzpicture}
\caption{Cliff Walking $4 \times 6$ grid. The optimal safe path (solid blue) goes Up from Start, Right along row~2, then Down to Goal.  Going Right from Start (dashed red) immediately falls off the cliff.}
\label{fig:cliffwalk-grid}
\end{figure}

\noindent The environment is defined arithmetically (avoiding Agda's literal-matching limits):

\begin{spec}
grid-row : ℕ → ℕ
grid-row s = if s <ᵇ 6 then 0 else if s <ᵇ 12 then 1 else if s <ᵇ 18 then 2 else 3

is-cliff s  = (grid-row s ≡ᵇ 3) ∧ ¬ᵇ (grid-col s ≡ᵇ 0) ∧ ¬ᵇ (grid-col s ≡ᵇ 5)
is-goal  s  = (grid-row s ≡ᵇ 3) ∧ (grid-col s ≡ᵇ 5)
reward-fn s = if is-goal s then 1 else 0
\end{spec}

\paragraph{Raw features.}
The same five-family feature taxonomy from FrozenLake applies: 4~\texttt{row-is} + 6~\texttt{col-is} + 4~\texttt{has-pos-reward} + 4~\texttt{is-self-loop} + 4~\texttt{leads-terminal} = 22 candidates, all 22 non-trivial.

\paragraph{Feature semantics and cliff detection.}
The dynamics features discover the cliff structure:

\begin{spec}
test-danger-13D : eval-cw (leads-terminal D) 13 ≡ true  = refl   -- row 2, col 1: D→19(cliff!)
test-danger-16D : eval-cw (leads-terminal D) 16 ≡ true  = refl   -- row 2, col 4: D→22(cliff!)
test-danger-12D : eval-cw (leads-terminal D) 12 ≡ false = refl   -- row 2, col 0: D→18(start, safe)
test-reward-17D : eval-cw (has-pos-reward D)  17 ≡ true = refl   -- row 2, col 5: D→23(goal, r=1)
\end{spec}

\noindent \texttt{leads-terminal~D} is \texttt{true} at row~2 columns~1--4 (the cells \emph{directly above the cliff}), but \texttt{false} at column~0 (above Start) and column~5 (above Goal).  The features automatically identify the ``danger zone'' as a band above the cliff, without any hand-coded notion of cliff proximity.

\paragraph{Optimal trajectory.}
The safe path takes 7~steps, avoids all cliff cells, and reaches the goal with reward~1:

\begin{spec}
test-trajectory : run-policy 18 10
  ≡ (U,12) ∷ (R,13) ∷ (R,14) ∷ (R,15) ∷ (R,16) ∷ (R,17) ∷ (D,23) ∷ []  = refl
test-rewards : collect-rewards 18 10 ≡ 0 ∷ 0 ∷ 0 ∷ 0 ∷ 0 ∷ 0 ∷ 1 ∷ []     = refl
test-all-safe : all-safe traj-states  ≡ true                                  = refl
\end{spec}

\noindent The FDMDP Finder independently verifies the optimal actions:

\begin{spec}
test-finder-18 : Finder.find-policy 18 8 ≡ U  = refl    -- Start: go Up
test-finder-12 : Finder.find-policy 12 8 ≡ R  = refl    -- Row 2, col 0: go Right
test-finder-17 : Finder.find-policy 17 8 ≡ D  = refl    -- Row 2, col 5: go Down to Goal
\end{spec}

\paragraph{CEGIS synthesis.}
With 22~features at depth~0, the version space has 24~candidates.  CEGIS synthesizes action-ranking predicates that encode spatial knowledge about the cliff:

\begin{spec}
obs-U≤D : (13 , false) ∷ (7 , true) ∷ []
synth-U≤D : synth-rank-pred 0 obs-U≤D ≡ just (feat (row-is 1))  = refl

obs-R≤D : (12 , true) ∷ (14 , false) ∷ []
synth-R≤D : synth-rank-pred 0 obs-R≤D ≡ just (feat (col-is 0))  = refl
\end{spec}

\noindent For ``is Down at least as good as Up?'', CEGIS discovers \texttt{row-is~1}: at row~1 it is safe to go Down (reaching row~2), while at row~2 going Down leads to the cliff.  For ``is Down at least as good as Right?'', it discovers \texttt{col-is~0}: only at the leftmost column (the Start corner) is Down~$\geq$~Right, since elsewhere in row~2 moving Down hits the cliff while Right stays safe.

\begin{remark}[Cliff Walking vs.\ FrozenLake]
While both are grid-based FDMDPs, they exercise the framework differently.  FrozenLake has \emph{scattered} terminal states (holes), whereas Cliff Walking has a \emph{contiguous line} of terminal states.  The discovered features reflect this: in FrozenLake, \texttt{leads-terminal} flags isolated danger spots (adjacent to scattered holes), while in Cliff Walking it identifies a \emph{band} of danger (the entire row above the cliff).  This demonstrates that the feature discovery pipeline adapts to different spatial topologies---a concrete benefit of the dynamics-based FTL.
\end{remark}


\subsection{Demo 9: Auto-FTL --- Zero Human Feature Engineering}
\label{sec:e2e-autoftl}

Demos~7 and~8 demonstrated fully automated policy synthesis, but both relied on \emph{human-defined} feature templates: \texttt{row-is}, \texttt{col-is}, and dynamics features specific to the grid interpretation.  A natural question is: can the features themselves be generated automatically?

We introduce the \emph{Automatic Feature Template Language} (Auto-FTL), a reusable module (\texttt{CSHRL.Synthesis.AutoFeatureNat}) that derives a complete feature set from only two inputs: the state count~$n$ and the step function.  No grid interpretation, no coordinate system, no domain knowledge.

\paragraph{Feature generation from state count.}
Given $n$~states encoded as $\{0, \ldots, n{-}1\}$, the Auto-FTL generates three feature families:

\begin{enumerate}
\item \textbf{State identity}: \texttt{state-is~$k$} for each $k \in \{0, \ldots, n{-}1\}$---one-hot per state.
\item \textbf{Factorization}: for each non-trivial divisor $w$ of $n$ (computed automatically via integer division):
  \begin{itemize}
  \item \texttt{mod-is~$w$~$k$}: is $s \bmod w = k$?
  \item \texttt{div-is~$w$~$k$}: is $\lfloor s / w \rfloor = k$?
  \end{itemize}
\item \textbf{Dynamics}: \texttt{has-pos-reward~$a$}, \texttt{is-self-loop~$a$}, \texttt{leads-terminal~$a$}---derived from the step function and terminal predicate.
\end{enumerate}

\noindent The arithmetic helpers (safe integer division, modulo, divisor enumeration) are defined with structurally decreasing fuel parameters, satisfying Agda's termination checker under \texttt{--safe}.

\paragraph{Emergent grid coordinates.}
For $n = 16$ (FrozenLake), the Auto-FTL discovers divisors $\{2, 4, 8\}$.  The factorization at $w = 4$ is particularly revealing:

\begin{spec}
test-auto-row3 : eval-auto (div-is 4 3) 14 ≡ true    = refl   -- div-is 4 k ≡ row-is k
test-auto-col2 : eval-auto (mod-is 4 2) 10 ≡ true     = refl   -- mod-is 4 k ≡ col-is k
\end{spec}

\noindent Since state $s$ in a 4-column grid satisfies $\text{row}(s) = \lfloor s/4 \rfloor$ and $\text{col}(s) = s \bmod 4$, the factorization features \emph{automatically recover the grid coordinates} without any grid interpretation.  The $w = 2$ and $w = 8$ factorizations additionally provide coarser/finer spatial partitions (even/odd, top/bottom half), giving CEGIS more options for discriminating states.

\paragraph{Feature enumeration.}
The complete enumeration for FrozenLake produces $16 + 28 + 12 = 56$ candidate features:

\begin{spec}
test-divisors       : divisors        ≡ 2 ∷ 4 ∷ 8 ∷ []     = refl
test-enum-count     : length enumerate-auto ≡ 56            = refl
test-discovered-count : length discovered   ≡ 56            = refl   -- all non-trivial
\end{spec}

\noindent All 56 features survive non-trivial filtering (none is constant across all 16 states).  This is $2.8\times$ the hand-crafted count of 20 features, but the version space remains tractable: 58 candidates at depth~0, approximately 6{,}800 at depth~1.

\paragraph{Fully automated pipeline with auto-features.}
The cascading CEGIS pipeline from Demo~7 works \emph{unchanged} with the auto-generated features:

\begin{spec}
pl = extract (synth-rank-pred 0 obs-is-L)          -- L-optimal predicate (depth 0)
pr = extract (synth-rank-pred 1 obs-is-R)          -- R-optimal predicate (depth 1)
auto-policy s = if eval pl s then L else if eval pr s then R else D
\end{spec}

\noindent CEGIS discovers predicates using the automatically generated features:
\begin{itemize}
\item \texttt{pl} uses a state-identity feature to capture the single L-optimal state.
\item \texttt{pr} combines a factorization feature (\texttt{div-is~4~3}, the emergent row coordinate) with a dynamics feature (\texttt{leads-terminal~D}) to capture R-optimal states---a conjunction of spatial and behavioral properties, discovered without human guidance.
\end{itemize}

\paragraph{Verification.}
The auto-policy is verified against the Finder at \emph{all} 11 non-terminal states:

\begin{spec}
test-auto-traj : run-auto 0 7
  ≡ (D,4) ∷ (D,8) ∷ (R,9) ∷ (D,13) ∷ (R,14) ∷ (R,15) ∷ []   = refl
test-auto-rewards    : auto-rewards 0 7 ≡ 0 ∷ 0 ∷ 0 ∷ 0 ∷ 0 ∷ 1 ∷ []  = refl
test-auto-safe       : all-safe auto-traj-states ≡ true                   = refl
test-auto-all-match  : check-all finder-map ≡ true                        = refl
\end{spec}

\noindent The trajectory, rewards, hole avoidance, and global Finder agreement are identical to the hand-crafted Demo~7.  The entire pipeline compiles in approximately 5 minutes.

\begin{remark}[From Hand-Crafted to Fully Automatic]
Demo~9 eliminates the last human input in the synthesis pipeline.  Demos~1--8 required human-defined feature templates---\texttt{row-is}/\texttt{col-is} for grids, \texttt{vals-eq}/\texttt{spread-eq} for placement problems.  The Auto-FTL replaces these with features derived from the state count's number-theoretic structure (divisors, modular arithmetic) and the step function's behavioral properties (reward, self-loops, terminal transitions).  The result is a synthesis pipeline where the human provides \emph{only} the environment definition ($n$, actions, step function), and every subsequent step---feature generation, feature filtering, observation generation, predicate synthesis, policy reconstruction, and correctness verification---is fully automatic and machine-verified.
\end{remark}


\subsection{Demo 10: Auto-FTL Cliff Walking --- Scaling to Richer Structure}
\label{sec:e2e-autoftl-cliffwalk}

Demo~9 applied the Auto-FTL to FrozenLake ($n = 16$, 3~divisors, 56~features).  We now apply it to Cliff Walking ($n = 24$, 6~divisors, 106~features)---a larger state space with richer factorization structure, a contiguous cliff, and a more complex optimal policy requiring three cascading CEGIS stages.  Again: zero human feature engineering, zero human observations.

\paragraph{Richer factorization from a highly composite number.}
The state count $n = 24$ has six non-trivial divisors, compared to three for $n = 16$:

\begin{spec}
test-divisors : divisors ≡ 2 ∷ 3 ∷ 4 ∷ 6 ∷ 8 ∷ 12 ∷ []   = refl
test-enum-count     : length enumerate-auto ≡ 106           = refl
test-discovered-count : length discovered   ≡ 106           = refl
\end{spec}

\noindent The 106 features comprise 24~state-identity, 70~factorization (across six divisors), and 12~dynamics features.  All survive non-trivial filtering.

\paragraph{Emergent grid coordinates for a non-square grid.}
For the $4 \times 6$ grid, the divisor $w = 6$ recovers the grid coordinates:

\begin{spec}
test-auto-row3 : eval-auto (div-is 6 3) 18 ≡ true    = refl   -- div-is 6 k ≡ row-is k
test-auto-col5 : eval-auto (mod-is 6 5) 17 ≡ true    = refl   -- mod-is 6 k ≡ col-is k
\end{spec}

\noindent This demonstrates that the Auto-FTL generalises beyond square grids.  The additional divisors ($w \in \{2, 3, 4, 8, 12\}$) provide complementary partitions: $w = 4$ groups states into blocks of 4, $w = 3$ into blocks of 3, and so on.  These coarser/finer partitions give CEGIS more discriminating power, which proves essential for the R-optimal predicate.

\paragraph{Three-stage cascading CEGIS.}
The Cliff Walking optimal policy at depth~7 assigns four distinct actions across the 19~non-terminal states: L~at state~0, U~at the start state~18, R~along row~2 (states~12--16), and D~everywhere else (states~1--11 and~17).  This requires three CEGIS stages instead of FrozenLake's two:

\begin{spec}
pl = extract (synth-rank-pred 0 obs-is-L)          -- L-optimal (depth 0)
pu = extract (synth-rank-pred 0 obs-is-U)          -- U-optimal (depth 0)
pr = extract (synth-rank-pred 1 obs-is-R)          -- R-optimal (depth 1)
auto-policy s = if eval pl s then L else if eval pu s then U
                else if eval pr s then R else D
\end{spec}

\noindent Stages~1 and~2 use depth-0 synthesis (singleton state identification).  Stage~3 requires depth~1: the R-optimal set $\{12, 13, 14, 15, 16\}$ cannot be captured by any single auto-feature (the closest, \texttt{div-is~6~2}, also includes state~17).  CEGIS discovers the disjunction \texttt{feat(state-is~16) $\vee$p feat(div-is~4~3)}, which combines a state-identity feature with a factorization feature: $\lfloor s / 4 \rfloor = 3$ covers states~12--15, and \texttt{state-is~16} adds the missing state.  This predicate is a depth-1 disjunction of atoms, discoverable within the version space.

\paragraph{Verification.}
The synthesised policy is verified against the Finder oracle at all 19~non-terminal states:

\begin{spec}
test-auto-traj : run-auto 18 10
  ≡ (U , 12) ∷ (R , 13) ∷ (R , 14) ∷ (R , 15) ∷
    (R , 16) ∷ (R , 17) ∷ (D , 23) ∷ []           = refl
test-auto-rewards  : auto-rewards 18 10 ≡ 0 ∷ 0 ∷ 0 ∷ 0 ∷ 0 ∷ 0 ∷ 1 ∷ []  = refl
test-auto-safe     : all-safe auto-traj-states ≡ true                         = refl
test-auto-all-match : check-all finder-map ≡ true                             = refl
\end{spec}

\noindent The optimal trajectory $18 \to 12 \to 13 \to 14 \to 15 \to 16 \to 17 \to 23$ navigates Up from Start, Right along row~2 (avoiding the cliff), then Down to Goal---collecting reward~1 with all intermediate states cliff-free.  The policy matches the Finder at every non-terminal state.

\begin{remark}[Scaling Auto-FTL]
This demo validates the Auto-FTL on a second, larger environment.  The key differences from FrozenLake are: (i)~the state count's richer divisor structure (6 vs.\ 3) provides more factorization features; (ii)~the non-square grid ($4 \times 6$) requires a different divisor ($w = 6$ vs.\ $w = 4$) for the grid coordinates, yet the Auto-FTL discovers this automatically; (iii)~the optimal policy spans four actions, requiring a three-stage cascade; (iv)~the R-optimal predicate requires depth-1 CEGIS, demonstrating that Auto-FTL features compose meaningfully under \texttt{PredProg}'s Boolean combinators.  The compilation takes approximately 4~hours---dominated by the Finder's depth-7 computation across 19~states without memoisation---confirming that the approach scales to environments with $\sim$20 non-terminal states at depth~7.
\end{remark}


\subsection{Demo 11: FrozenLake \texorpdfstring{$8 \times 8$}{8x8} --- Scaling to 64 States}
\label{sec:e2e-frozenlake8x8}

Demo~10 applied the Auto-FTL to Cliff Walking ($n = 24$, 19~non-terminal states).  We now scale to FrozenLake~$8{\times}8$ ($n = 64$, 53~non-terminal states, 10~holes)---a $4\times$ increase in state space over the $4{\times}4$ variant.  This demo introduces two infrastructure improvements required for scalable synthesis: a \emph{memoized Finder} that computes optimal policies in $O(d \cdot n^2)$ time via dynamic programming, and a \emph{greedy cascading CEGIS} (\texttt{synth-greedy-or}) that handles scattered state sets unreachable by bounded-depth predicate enumeration.

\paragraph{The environment.}
The standard OpenAI Gym $8{\times}8$ FrozenLake map has 64 states: Start at~0, Goal at~63, and 10 holes at positions $\{19, 29, 35, 41, 42, 46, 49, 52, 54, 59\}$.  Movement is defined arithmetically using \texttt{divℕ}/\texttt{modℕ} for row/column extraction---no 64-case pattern match:

\begin{spec}
go r c L = if c ≡ᵇ 0 then r * 8 + c else r * 8 + (c ∸ 1)
go r c D = if r ≡ᵇ 7 then r * 8 + c else (r + 1) * 8 + c
go r c R = if c ≡ᵇ 7 then r * 8 + c else r * 8 + (c + 1)
go r c U = if r ≡ᵇ 0 then r * 8 + c else (r ∸ 1) * 8 + c
move s a = if is-terminal s then s else go (divℕ s 8) (modℕ s 8) a
\end{spec}

\paragraph{Memoized Finder.}
The original Finder evaluates the full state-action tree, which is exponential in the planning depth.  For 64 states at depth~16, this is computationally infeasible.  The memoized Finder replaces tree exploration with \emph{bottom-up dynamic programming}: starting from a base layer of empty traces, each layer computes the optimal trace at every state by looking up successor values in the previous layer.  The DP table \texttt{dp~$d$} has complexity $O(d \cdot n \cdot |A|)$, making depth-16 computation over 64 states tractable.

A key Agda-specific optimization: the \texttt{Memoized} module exports a \texttt{policy-table : List~Action} that computes all 64 optimal actions in a \emph{single} pass over the DP table.  This avoids Agda's normalizer re-evaluating the DP table for each state lookup, reducing 53 independent evaluations to one shared computation.

\paragraph{Auto-FTL features.}
The state count $n = 64$ has five non-trivial divisors:

\begin{spec}
test-divisors : divisors ≡ 2 ∷ 4 ∷ 8 ∷ 16 ∷ 32 ∷ []    = refl
\end{spec}

\noindent The divisor $w = 8$ recovers the $8{\times}8$ grid coordinates: \texttt{div-is~8~$k$} $\equiv$ \texttt{row-is~$k$}, \texttt{mod-is~8~$k$} $\equiv$ \texttt{col-is~$k$}.  The 200 discovered features comprise 64~state-identity, 120~factorization (across five divisors), and 16~dynamics features.

\paragraph{Optimal policy structure.}
The Finder at depth~16 partitions the 53 non-terminal states as: D-optimal $= 33$ states, R-optimal $= 17$ states, U-optimal $= \{34, 51, 58\}$ (3~states), L-optimal $= 0$ states.  The U-optimal states correspond to grid positions $(4,2)$, $(6,3)$, $(7,2)$---positions where going Right or Down leads into holes, and going Up is the only safe direction toward the goal.

\paragraph{Greedy cascading CEGIS.}
The 3 U-optimal states $\{34, 51, 58\}$ are \emph{scattered}: no single auto-feature or depth-1 combination captures exactly this set.  The standard \texttt{synth-rank-pred} requires a single \texttt{PredProg} at bounded depth that \emph{exactly} classifies all observations, which fails for scattered sets.  We introduce \texttt{synth-greedy-or}, which builds a disjunction of atomic predicates iteratively:
\begin{enumerate}
\item Find an atom with \emph{zero false positives} and at least one true positive.
\item Remove covered true observations.
\item Recurse, combining results with $\vee$p.
\end{enumerate}

\noindent For the U-optimal set, this produces \texttt{feat(state-is~34)~$\vee$p~feat(state-is~51)~$\vee$p~feat(state-is~58)~$\vee$p~falsep}---three state-identity features chained by disjunction.  Each step searches only the 202 atoms (not the exponential depth-2+ space), keeping synthesis tractable.  For the 17 R-optimal states, \texttt{synth-greedy-or} may find broader features (factorization or dynamics) that cover multiple states per step, producing a more compact predicate.

\paragraph{Bundled pipeline.}
The entire synthesis pipeline---DP table computation, finder-map derivation, greedy CEGIS for U and R, and verification---runs in a single \texttt{let}-chain.  This ensures the DP table is normalized exactly once by Agda's evaluator:

\begin{spec}
pipeline : List Action → Bool
pipeline ptbl =
  let fm    = map (λ s → (s , nth-act ptbl s)) non-terminal-states
      obs-U = map (λ { (s , a) → (s , a ≟ᵃ U) }) fm
      pu    = synth-greedy-or 10 obs-U
      rem-U = bfilter-pair (λ { (s , a) → not (a ≟ᵃ U) }) fm
      obs-R = map (λ { (s , a) → (s , a ≟ᵃ R) }) rem-U
      pr    = synth-greedy-or 20 obs-R
  in check-with pu pr fm
test-pipeline : pipeline Memo.policy-table ≡ true    = refl
\end{spec}

\paragraph{Verification and performance.}
The synthesized policy matches the Finder at all 53 non-terminal states, proved by \texttt{refl}.  The full file---environment definition, Auto-FTL instantiation, memoized Finder, greedy CEGIS, and verification---compiles in approximately \textbf{7~seconds}.  Compare this to the unmemoized Cliff Walking demo ($n = 24$, depth~7), which requires $\sim$4~hours: the memoized Finder provides a ${\sim}2{,}000\times$ speedup, making depth-16 computation over 64 states faster than depth-7 over 24 states.

\begin{remark}[Scaling Bottlenecks and Solutions]
This demo exposed two scalability bottlenecks and their solutions.  First, Agda's normalizer recomputes top-level definitions independently for each reference: the DP table (thousands of nested \texttt{suc} constructors) was re-evaluated once per state lookup.  The fix---\texttt{policy-table} compresses the DP output into a compact \texttt{List~Action} (64 constructors), which is cheap to copy.  Second, the standard \texttt{synth-rank-pred} enumerates all PredProg terms up to a depth bound, which grows as $O(n^{2^d})$.  With 200~features, depth~2 produces $\sim$8~billion candidates---infeasible.  The fix---\texttt{synth-greedy-or}---avoids depth enumeration entirely, building the predicate incrementally from atoms.  These two optimizations transform a multi-hour computation into a 7-second one.
\end{remark}


\subsection{Demo 12: Taxi \texorpdfstring{$5 \times 5$}{5x5} --- Relational Navigation at Scale}
\label{sec:e2e-taxi}

Demos~9--11 operated on environments with at most 64~states and a single reward structure (hole avoidance or cliff avoidance).  We now tackle Taxi~$5{\times}5$---the classic OpenAI Gym benchmark~\cite{dietterich2000hierarchical}---with \textbf{501~states}, \textbf{6~actions}, \textbf{internal walls}, and a multi-phase task (navigate to passenger, pick up, navigate to destination, drop off).  This represents a ${\sim}10\times$ scale-up over FrozenLake~$8{\times}8$ and the first demo requiring \emph{relational} reasoning about target-relative position.

\paragraph{Environment.}
The Taxi environment is a $5 \times 5$ grid with four designated locations: R$(0{,}0)$, G$(0{,}4)$, Y$(4{,}0)$, B$(4{,}3)$.  Internal walls block East/West movement: between columns~1--2 at rows~0--1, and between columns~0--1 and~2--3 at rows~3--4.  A state encodes $(\text{row}, \text{col}, \text{passenger}, \text{dest})$ as a single natural number: $\text{row} \times 100 + \text{col} \times 20 + \text{pass} \times 4 + \text{dest}$, yielding 500~regular states plus a terminal state~500.  There are 6~actions: South, North, East, West, Pickup, and Dropoff.  Reward is~1 at successful dropoff, 0~otherwise.

\paragraph{Memoized Finder at scale.}
The memoized Finder is instantiated at depth~25 over 501~states.  Despite the $10\times$ state-space increase, the DP table computation completes in approximately \textbf{60~seconds} (CPU time)---demonstrating that the memoized infrastructure scales well to hundreds of states.  The maximum optimal trajectory length is 22~steps (e.g., navigating diagonally across the grid around walls, picking up, navigating to the opposite corner, and dropping off).

\paragraph{Action distribution.}
The optimal policy assigns 6~distinct actions across the 500~non-terminal states: S${}=180$, N${}=220$, E${}=35$, W${}=45$, P${}=16$, X${}=4$.  The dominance of N over S reflects the Finder's tie-breaking (N precedes E/W in the action list, so states where vertical and lateral routes are equally short prefer N).  The 16~Pickup states correspond to~$4$~locations~$\times$~$4$~destinations, and the~4 Dropoff states to the~4 at-destination-with-passenger-aboard configurations.

\paragraph{Compound features for Pickup and Dropoff.}
The Auto-FTL generates ${\sim}1{,}700$~features for $n = 500$ (10~divisors including 4, 20, and~100, which decompose the encoding perfectly).  We define two targeted \emph{compound features}:
\begin{itemize}
\item \texttt{at-dest-aboard}: passenger in taxi \textbf{and} taxi at destination (covers X-optimal states).
\item \texttt{at-pass-waiting}: passenger not in taxi \textbf{and} taxi at passenger's location (covers P-optimal states).
\end{itemize}
\noindent Both are verified against the Finder at all 500~states.

\paragraph{Relational wall-aware navigation.}
For interpretability, we also provide a hand-crafted navigation synthesis.  Unlike previous demos where feature-based predicates sufficed, Taxi navigation requires \emph{relational reasoning}: the optimal action depends on the taxi's position \emph{relative to the target}, which itself depends on whether the passenger is aboard (target~$=$~destination) or waiting (target~$=$~passenger location).  Furthermore, the internal walls create \emph{detour corridors}: row~2 is wall-free and serves as a ``highway'' for lateral movement.

The synthesized navigation function \texttt{nav} encodes three rules:
\begin{enumerate}
\item \textbf{Column matches target:} go vertically (S if below, N if above).
\item \textbf{Row matches target:} go laterally (E or W), unless walls block the path, in which case detour through row~2.
\item \textbf{Both differ:}
  \begin{itemize}
  \item At row~2 (highway): go laterally if the \emph{next} row toward the target has a wall blocking the lateral path; otherwise go vertically (wins tie-break).
  \item Adjacent to target row (rows $\{1{,}0\}$ or $\{3{,}4\}$) with wall at target row: detour toward row~2.
  \item Otherwise: go vertically (S/N wins tie-break over E/W when paths are equally short).
  \end{itemize}
\end{enumerate}

\noindent The \texttt{blocked-east} and \texttt{blocked-west} functions check whether the \emph{entire lateral path} from the current column to the target column is obstructed by any wall in the given row---not just the immediate cell.  This global path check, combined with the row-2 highway insight, produces a navigation function that matches the Finder's globally optimal choices at \textbf{all 500~states}.

\paragraph{Interpretable pipeline and verification.}
The complete synthesized policy composes the three components:

\begin{spec}
synth-policy s =
  if at-dest-aboard s then X
  else if at-pass-waiting s then P
  else nav (row s) (col s) (target-row s) (target-col s)

test-pipeline : pipeline Memo.policy-table ≡ true
test-pipeline = refl
\end{spec}

\noindent The \texttt{refl} proof verifies that \texttt{synth-policy} produces the same action as the Finder's optimal policy for every non-terminal state.

\paragraph{Resolving the Peano Bottleneck: fully automated CEGIS.}
Initial experiments revealed that custom recursive \texttt{divℕ}/\texttt{modℕ} functions imposed $O(n)$ reduction cost per arithmetic operation, making iterative CEGIS infeasible at 500~states.  However, Agda's standard library provides \texttt{\_/\_} and \texttt{\_\%\_}, which delegate to the \texttt{BUILTIN NATDIVSUCAUX}/\texttt{NATMODSUCAUX} pragmas---backed by GMP arbitrary-precision arithmetic and evaluated in $O(1)$ time by the normalizer.  Replacing the custom functions with these builtin-backed wrappers resolved the bottleneck entirely.

With $O(1)$ arithmetic, the same Auto-FTL \texttt{discovered} features (${\sim}1{,}700$ features) can drive a \textbf{fully automated 6-stage cascading CEGIS} pipeline:

\begin{spec}
auto-pipeline ptbl =
  let fm   = map (λ s → (s , nth-act ptbl s)) non-terminal-states
      pX   = synth-greedy-or 10  (obs-for X fm)   -- 4 states
      pP   = synth-greedy-or 20  (obs-for P ...)   -- 16 states
      pS   = synth-greedy-or 200 (obs-for S ...)   -- 180 states
      pN   = synth-greedy-or 250 (obs-for N ...)   -- 220 states
      pE   = synth-greedy-or 50  (obs-for E ...)   -- 35 states
                                                    -- W default: 45
  in check-auto pX pP pS pN pE fm

test-auto-pipeline : auto-pipeline Memo.policy-table ≡ true
test-auto-pipeline = refl
\end{spec}

\noindent This pipeline requires \textbf{zero human-provided features}: the Auto-FTL generates features from the state count and step function, the memoized Finder computes the optimal policy table, observations are derived automatically, and greedy cascading CEGIS synthesizes a \texttt{PredProg} disjunction per action.  The entire file---environment definition, memoized Finder, both pipelines, and verification---compiles in approximately \textbf{9.5~minutes}.

\begin{remark}[BUILTIN Arithmetic and the Peano Bottleneck]
\label{rem:builtin-arith}
The resolution of the Peano Bottleneck reveals an important architectural insight.  Agda's natural numbers are \emph{externally} Peano-style (constructors \texttt{zero} and \texttt{suc}), but \emph{internally} backed by GMP big integers when the \texttt{BUILTIN NATURAL} pragma is in effect.  Operations that pattern-match on successor structure (\texttt{suc}-recursive division, modulo) force Peano unfolding at $O(n)$; operations routed through \texttt{BUILTIN} pragmas (\texttt{\_/\_}, \texttt{\_\%\_}, \texttt{\_==\_}) bypass the unary representation entirely.  The fix was surgical: replacing 10~lines of custom arithmetic with 4~lines of stdlib wrappers, yielding a $2.5\times$ speedup for the hand-crafted pipeline and, more critically, making full automated CEGIS tractable at 500~states.
\end{remark}

\begin{center}
\small
\begin{tabular}{@{}lccccc@{}}
& \textbf{Cliff Walking} & \textbf{FL $4{\times}4$} & \textbf{FL $8{\times}8$} & \textbf{Taxi $5{\times}5$} \\
\hline
States (non-terminal) & 24 (19) & 16 (11) & 64 (53) & 501 (500) \\
Actions & 4 & 4 & 4 & 6 \\
Features (auto) & 106 & 56 & 200 & ${\sim}1{,}700$ \\
Finder depth & 7 & --- & 16 & 25 \\
Memoized & no & --- & yes & yes \\
Policy stages & 3 & 2 & 2 & 6 (auto) / 3 (interp.) \\
Synthesis & CEGIS & CEGIS & CEGIS & greedy-or / decision-tree \\
Total compile time & ${\sim}$4 hours & $<$25 s & ${\sim}$7 s & ${\sim}$13 min (DT) \\
\end{tabular}
\end{center}

\noindent Taxi introduces four qualitative advances over the previous demos: (1)~\emph{relational} predicates that reason about position relative to a dynamically determined target, (2)~a \emph{wall-aware} navigation strategy that exploits spatial structure (the row-2 highway), (3)~the first demonstration that \textbf{fully automated CEGIS scales to 500~states and 6~actions} once arithmetic bottlenecks are resolved, and (4)~a \textbf{greedy decision-tree synthesizer} that produces interpretable policies from relational features.

\paragraph{Feature expressiveness experiments.}
We investigated four feature sets to characterize the expressiveness boundary:

\begin{enumerate}
\item \textbf{Compact positional features} (\texttt{discovered-compact}, 2{,}851 features): drops the 500~\texttt{state-is} atoms, retaining factorization, threshold, and dynamics features.  \textbf{Fails} for navigation stages: no single modular, divisor, or threshold feature can separate ``go south'' from ``go north'' states, because the same absolute row can be above or below the target depending on the passenger component.

\item \textbf{Relational features with flat disjunctions} (12 features): target-relative comparisons (\texttt{row~<~target-row}, \texttt{col~=~target-col}), wall-blocked predicates, and \texttt{pass-aboard}/\texttt{at-highway} flags.  These are the right \emph{vocabulary}---they capture the comparisons the hand-crafted \texttt{nav} function uses.  Depth-1 greedy CEGIS (disjunctions of pairwise conjunctions, ${\sim}456$~atoms) \textbf{fails}: Taxi's navigation logic is a depth-4+ decision tree, and no disjunction of depth-1 conjunctions can express this structure.  The problem is the \emph{algorithm}, not the \emph{features}.

\item \textbf{Lookup-table CEGIS} (\texttt{discovered} with \texttt{state-is}, ${\sim}1{,}700$ features): \textbf{succeeds} in 9.5~minutes, but produces opaque 180-term disjunctions of state-identity atoms.

\item \textbf{Decision-tree CEGIS with relational features} (14 features): the same relational vocabulary augmented with two absolute-position features (\texttt{row-above-hwy}, \texttt{trow-above-hwy}) needed for detour-direction decisions, synthesized using a new \texttt{synth-decision-tree} algorithm.  \textbf{Succeeds} in ${\sim}13$~minutes, producing interpretable if-then-else trees verified optimal at all 500~states.
\end{enumerate}

\paragraph{Decision-tree synthesis.}
The \texttt{synth-decision-tree} algorithm addresses the structural limitation of \texttt{synth-greedy-or}.  While the latter builds flat disjunctions (each atom must have zero false positives), the decision-tree synthesizer greedily splits observations on the feature that minimizes Gini impurity ($\text{TP}\times\text{FP} + \text{FN}\times\text{TN}$), then recurses on each branch.  The resulting \texttt{PredProg} encodes an if-then-else tree: splitting on feature~$f$ yields $(f \wedge p_{\text{true}}) \vee (\neg f \wedge p_{\text{false}})$, which evaluates to $p_{\text{true}}$ when $f$ holds and $p_{\text{false}}$ otherwise.  The algorithm takes a fuel parameter for termination; with 14~relational atoms and fuel~14, it produces trees that correctly classify all navigation states.

The critical insight from the expressiveness experiments is that the 12 original relational features are provably insufficient: states at $(0, 3, \text{trow}{=}0, \text{tcol}{=}0)$ and $(4, 3, \text{trow}{=}4, \text{tcol}{=}0)$ have identical values for all 12~features (same relative comparisons, same wall patterns at the grid extremes) but require opposite actions (S~vs.~N for highway detour direction).  Adding \texttt{row-above-hwy} ($\text{row} < 2$) and \texttt{trow-above-hwy} ($\text{target-row} < 2$) resolves this ambiguity.

\noindent The four-way comparison reveals the full \emph{interpretability landscape}: (1)~positional features lack relational expressiveness; (2)~relational features with flat disjunctions have the right vocabulary but the wrong algorithm; (3)~lookup-table CEGIS succeeds but produces opaque results; (4)~\textbf{decision-tree CEGIS with relational features achieves both correctness and interpretability}---the synthesized trees mirror the hand-crafted \texttt{nav} function's logic while being fully automatically generated.


\subsection{Analysis: Sample Efficiency and Policy Compression}
\label{sec:analysis-compression}

The fully automated demos reveal two striking properties: CEGIS converges with far fewer observations than the Finder produces, and the synthesised policies are orders of magnitude smaller than their deep RL counterparts.

\paragraph{Finder vs.\ CEGIS: exhaustive oracle, frugal learner.}
The Finder is an \emph{exhaustive} oracle: for each non-terminal state, it evaluates every action sequence to the depth horizon, computing the full lexicographic trace.  This brute-force computation dominates compilation time (depth~7 across 19~states requires $\sim$4~hours without memoisation).  CEGIS, by contrast, is a \emph{frugal learner}: each observation $(s, b)$ eliminates all version-space candidates $p$ where $\texttt{eval}\;p\;s \neq b$.  Version-space shrinkage is monotonic (\texttt{vs-shrinks}), and the Propagation Theorem (\texttt{refine-equiv}) guarantees that feature-equivalent states yield \emph{definitionally identical} refinement---observing at one subsumes observing at the other.

The cascading pipeline processes 19~states per stage across 3~stages (L, U, R), totalling \textbf{56 oracle queries}.  Each query labels one state as positive or negative for the target action.  The table below shows VS convergence verified by Agda \texttt{refl}:

\begin{center}
\small
\begin{tabular}{@{}lccc@{}}
\textbf{Stage} & \textbf{Initial VS} & \textbf{Final VS} & \textbf{Reduction} \\
\hline
L-optimal (depth 0) & 108 & 15 & 86\% \\
U-optimal (depth 0) & 108 & 6 & 94\% \\
R-optimal (depth 1) & 23{,}544 & 2 & 99.99\% \\
\end{tabular}
\end{center}

\noindent In general, convergence at depth~$d$ with $n$~features requires $O(\log n^{2^d})$ observations in the best case---logarithmic in the version-space size.  The 19 Finder queries per stage are generous; CEGIS converges well before exhausting them.  The full Finder map is computed for \emph{verification} (\texttt{test-auto-all-match}), not because CEGIS needs it.


\paragraph{Learning curve: oracle queries vs.\ goal success rate.}
To present the results in the standard RL format, Figure~\ref{fig:learning-curve} plots policy performance against cumulative oracle queries.  The Finder partitions the 19 non-terminal states as: D-optimal $= \{0\text{--}11, 17\}$ (13 states), R-optimal $= \{12\text{--}16\}$ (5 states), U-optimal $= \{18\}$ (1 state).  At each stage boundary, we evaluate the partial policy's goal success rate (fraction of starting states from which the goal is reachable):

\begin{center}
\small
\begin{tabular}{@{}rlcc@{}}
\textbf{Queries} & \textbf{Policy} & \textbf{Accuracy} & \textbf{Success} \\
\hline
0  & D everywhere           & 13/19 (68\%) & 3/19 (16\%) \\
38 & U at state 18, else D  & 14/19 (74\%) & 3/19 (16\%) \\
56 & full auto-policy       & 19/19 (100\%) & 19/19 (100\%) \\
\end{tabular}
\end{center}

\noindent The accuracy and success values are derived from the verified Finder output (\texttt{finder-map}) and verified auto-policy (\texttt{test-auto-all-match = refl}).

\begin{figure}[H]
\centering
\begin{tikzpicture}
\begin{axis}[
    width=0.92\linewidth,
    height=5.8cm,
    xlabel={Training samples (log scale)},
    ylabel={Goal success rate (\%)},
    xmode=log,
    xmin=1, xmax=15000,
    ymin=0, ymax=110,
    ytick={0,20,40,60,80,100},
    legend style={at={(0.98,0.42)}, anchor=east, font=\small,
                  draw=nord3, fill=nord6, fill opacity=0.9},
    grid=major,
    grid style={nord4, thin},
    axis line style={nord3},
    tick style={nord3},
    every axis label/.style={font=\small},
    every tick label/.style={font=\footnotesize},
]

% DQN schematic (typical published results for Cliff Walking)
\addplot[nord11, thick, dashed, smooth, mark=none]
    coordinates {
        (1,0) (50,5) (200,20) (500,40)
        (1000,60) (2000,75) (3500,85)
        (5000,90) (8000,93) (12000,95)
    };
\addlegendentry{DQN ($\sim$10K ep., $\approx$95\%)}

% VIPER (Bastani et al. 2018): DRL training + decision tree extraction
\addplot[nord12, thick, dashdotted, smooth, mark=none]
    coordinates {
        (1,0) (50,3) (200,15) (500,35)
        (1000,55) (2000,72) (3500,84)
        (5000,92) (7000,98) (8000,100)
    };
\addlegendentry{VIPER ($\sim$8K ep.+extract, 100\%)}

% CSH synthesis (verified data points)
\addplot[nord10, ultra thick, mark=*, mark size=3pt]
    coordinates {
        (1,15.8) (38,15.8) (56,100)
    };
\addlegendentry{CSH (56 queries, \textbf{100\% proven})}

% Annotation: "proven optimal"
\node[anchor=south west, font=\footnotesize\bfseries, nord10]
    at (axis cs:60,100) {proven};

% Annotation: "no proof"
\node[anchor=south, font=\footnotesize, nord11]
    at (axis cs:12000,95) {no proof};

% Annotation for VIPER
\node[anchor=north east, font=\footnotesize, nord12]
    at (axis cs:7000,96) {post-hoc};

% Vertical line at 56
\draw[nord10, thin, densely dotted]
    (axis cs:56,0) -- (axis cs:56,100);

\end{axis}
\end{tikzpicture}
\caption{Goal success rate vs.\ training samples (log scale).  \textbf{CSH synthesis} (solid blue) achieves provably 100\% optimal performance in 56 oracle queries.  \textbf{VIPER}~\cite{Bastani2018} (dash-dotted orange) illustrates the distillation approach: a DNN is trained via standard RL, then compressed into a decision tree and verified post-hoc for specific properties---the total sample cost includes both training and extraction.  A schematic \textbf{DQN} (dashed red) shows typical deep RL convergence.  CSH's step-like curve reflects the fundamental difference: synthesis produces a \emph{complete, proven} policy once sufficient observations accumulate, rather than gradually approximating one.  VIPER can achieve perfect performance but cannot guarantee optimality---it verifies only specific properties (robustness, stability), and counterexamples in the extracted tree are discovered post-hoc~\cite{Bastani2018}.}
\label{fig:learning-curve}
\end{figure}

\paragraph{Policy compression: 38~bits vs.\ 192{,}640~bits.}
The synthesised Cliff Walking policy is a cascading \texttt{if-then-else} over three \texttt{PredProg} terms:

\begin{center}
\small
\begin{tabular}{@{}ll@{}}
\texttt{falsep} & $\to$ L \quad\emph{(dead---no L-optimal states at this horizon)} \\
\texttt{state-is~18} & $\to$ U \\
\texttt{(state-is~16) $\vee$p (div-is~4~3)} & $\to$ R \\
\texttt{else} & $\to$ D \\
\end{tabular}
\end{center}

\noindent The first branch is dead code: at the 8-step planning horizon, no state is L-optimal, so CEGIS synthesises \texttt{falsep}.  The effective policy uses only two predicates.  Encoding the full 4-stage structure requires: 3~feature indices from 106 ($3 \times \lceil \log_2 106 \rceil = 21$~bits), 3~Boolean structure tags ($\sim$5~bits), 4~action labels ($4 \times 2 = 8$~bits), plus overhead ($\sim$4~bits): approximately \textbf{38~bits} total.

A standard DQN for the same environment (24-input one-hot, two 64-unit hidden layers, 4-output) has 6{,}020~float32 parameters: \textbf{192{,}640~bits ($\sim$24\,KB)}.  Even quantised to int8, a DQN requires 48{,}160~bits ($\sim$6\,KB).

\begin{center}
\small
\begin{tabular}{@{}lccc@{}}
& \textbf{CSH Synthesis} & \textbf{VIPER}~\cite{Bastani2018} & \textbf{DQN} \\
\hline
Policy size & $\sim$38 bits & 31--9{,}757 nodes & $\sim$192{,}640 bits \\
Compression & \textbf{1$\times$} & $\sim$50$\times$ larger & $\sim$5{,}000$\times$ larger \\
Inference cost & 3 int.\ comparisons & tree traversal & $\sim$6{,}000 mul-adds \\
Correctness & Proven (\texttt{refl}) & Post-hoc (Z3/SOS) & Approximate \\
Interpretability & Human-readable & Partially readable & Opaque \\
Training & 56 oracle queries & ${\sim}$10K ep.\ + extract & $\sim$10{,}000+ ep. \\
Hyperparameters & 0 & ${\sim}$10 & $\sim$8 \\
Counterexamples & Impossible & Discovered post-hoc & Unknown \\
\end{tabular}
\end{center}

\paragraph{Why the gap?  Kolmogorov complexity.}
The optimal Cliff Walking policy has very low Kolmogorov complexity: the environment's regular spatial structure (grid, contiguous cliff, single goal) induces a policy describable by a few spatial predicates.  The Auto-FTL \emph{discovers} this regularity---grid coordinates emerge from modular arithmetic---and \texttt{PredProg} \emph{compresses} it into a near-minimal Boolean program.  A DQN, by contrast, is a generic function approximator: it must represent this inherently simple mapping using thousands of floating-point parameters, ``wasting'' capacity on a structure that admits a 38-bit description.

VIPER's decision trees sit between these extremes: they are more structured than DNNs (axis-aligned splits) but less compact than Boolean predicate programs (each node repeats a state-space partition).  A VIPER tree for a similar grid world would typically require 31--100+~nodes, each encoding a threshold comparison---roughly 10--50$\times$ larger than our \texttt{PredProg}.  Note that our \texttt{synth-decision-tree} (Demo~12) also builds decision trees, but within the \texttt{PredProg} framework: each tree node is a relational feature, and the result is a verified \texttt{PredProg} term.  Unlike VIPER, which distills an unverified DNN, our decision trees are synthesized directly from the Finder's proven-optimal policy and verified at every state.

This compression is not an artefact of small environments.  For \emph{any} environment whose optimal policy has low K-complexity (regular structure, spatial patterns, symmetries), the CSH pipeline will produce a proportionally compact \texttt{PredProg}---while a DQN's parameter count is determined by architecture, not by the policy's intrinsic complexity.  The FTL controls this: by providing features aligned with the environment's structure, it ensures the optimal predicate lies at low \texttt{PredProg} depth, where the encoding is minimal.

\begin{remark}[Beyond compression]
The 5{,}000$\times$ size advantage over DQN understates the practical gap.  The synthesised policy is \emph{verified}: the \texttt{test-auto-all-match} proof, checked by Agda's kernel, certifies optimality at every state.  No amount of DQN training or testing can provide this guarantee---and VIPER's post-hoc verification checks only specific properties (robustness, stability), not global optimality.  Furthermore, the synthesised predicates are \emph{explanatory}: ``go Right when $\lfloor s/4 \rfloor = 3$ or $s = 16$'' tells a human \emph{why} the action is optimal, relating it to the state's position in the factorization structure.  These properties---provable optimality, interpretability, minimal encoding---flow directly from the CSH foundation: the Finder computes the ground truth (CSH-backed optimality), Auto-FTL provides the language (structure discovery), and CEGIS compresses it (sample-efficient synthesis).
\end{remark}

\section{Cross-Pair Propagation: Learning A~$>$~B from C~vs~D}
\label{sec:cross-pair}

We now demonstrate the deepest consequence of the Propagation Theorem: CEGIS observations are \emph{pair-agnostic}.  When ranking pairs share the same feature structure, observing a sub-optimal pair determines the top-tier ranking---without ever directly comparing the actions of interest.

\subsection{Setup: 4 Actions, 2 Features, 6 States}

Consider a scenario with 4 actions (A, B ``top tier''; C, D ``bottom tier''), 2 Boolean features, and 6 states forming 3 equivalence classes:

\begin{spec}
data State : Set where s₁ s₂ s₃ s₄ s₅ s₆ : State
data Action : Set where A B C D : Action
data Feature : Set where f₁ f₂ : Feature

eval-feat : Feature → State → Bool
-- Class α (s₁, s₂): f₁ = T, f₂ = T
-- Class β (s₃, s₄): f₁ = F, f₂ = T
-- Class γ (s₅, s₆): f₁ = T, f₂ = F
\end{spec}

\noindent The action ranking is entirely determined by feature $f_1$:
\begin{itemize}
\item $f_1 = \mathsf{true}$: $A > B > C > D$
\item $f_1 = \mathsf{false}$: $D > C > B > A$
\end{itemize}
So \texttt{prefer(A,B)} = \texttt{prefer(C,D)} = \texttt{feat~f₁}.  All 6 ranking pairs are governed by the \emph{same} predicate.

\subsection{CEGIS from C-vs-D Observations Only}

We feed CEGIS \emph{only} C-vs-D comparisons and track the version space:

\begin{spec}
vs₀ : VersionSpace
vs₀ = initial-vs 0

check-vs₀ : length vs₀ ≡ 4       -- truep, falsep, feat f₁, feat f₂
check-vs₀ = refl

obs₁ : PredObs                    -- C > D at s₁ (class α)
obs₁ = s₁ , true

check-vs₁ : length (refine vs₀ obs₁) ≡ 3   -- falsep eliminated
check-vs₁ = refl

obs₂ : PredObs                    -- D > C at s₃ (class β)
obs₂ = s₃ , false

check-vs₂ : length (refine (refine vs₀ obs₁) obs₂) ≡ 1
check-vs₂ = refl                  -- only feat f₁ survives!

pinpointed : cegis-loop vs₀ (obs₁ ∷ obs₂ ∷ []) ≡ feat f₁ ∷ []
pinpointed = refl
\end{spec}

\noindent Two C-vs-D observations pinpoint \texttt{feat~f₁} as the sole survivor.  The evolution:

\medskip
\begin{center}
\begin{tabular}{llcc}
\textbf{Step} & \textbf{Observation} & \textbf{|VS|} & \textbf{Eliminated} \\
\hline
Initial & --- & 4 & --- \\
1 & $C > D$ at $s_1$ (class $\alpha$) & 3 & \texttt{falsep} \\
2 & $D > C$ at $s_3$ (class $\beta$) & 1 & \texttt{truep}, \texttt{feat~f₂} \\
\end{tabular}
\end{center}

\subsection{The Payoff: A~$>$~B Without Observing A or B}

We \emph{never} observed A vs B.  Yet the surviving program correctly determines the A-vs-B ranking at all 6 states:

\begin{spec}
ab-s₁ : eval (feat f₁) s₁ ≡ true    -- A > B at class α
ab-s₁ = refl

ab-s₃ : eval (feat f₁) s₃ ≡ false   -- B > A at class β
ab-s₃ = refl

ab-s₅ : eval (feat f₁) s₅ ≡ true    -- A > B at class γ
ab-s₅ = refl
\end{spec}

\noindent In fact, \emph{all 6 ranking pairs} are determined: $\texttt{prefer}(X,Y) = \texttt{feat}~f_1$ for every $(X,Y) \in \{(A,B), (A,C), (A,D), (B,C), (B,D), (C,D)\}$.  Two sub-optimal observations determined 6 rankings at 6 states.

\subsection{Dissolution Multiplies the Savings}

Feature-equivalent states add zero information:

\begin{spec}
s₁≈s₂ : AllFeatAgree s₁ s₂
s₁≈s₂ f₁ = refl
s₁≈s₂ f₂ = refl

dissolution-α : refine (refine vs₀ (s₁ , true)) (s₂ , true)
              ≡ refine vs₀ (s₁ , true)
dissolution-α = refine-absorb vs₀ s₁ s₂ true s₁≈s₂
\end{spec}

\noindent The total savings:

\medskip
\begin{center}
\begin{tabular}{lcc}
\textbf{Approach} & \textbf{Observations} & \textbf{Factor} \\
\hline
Na\"ive (all pairs $\times$ all states) & $6 \times 6 = 36$ & $1\times$ \\
Dissolution only (all pairs $\times$ classes) & $6 \times 3 = 18$ & $2\times$ \\
Cross-pair + dissolution & $2$ & $18\times$ \\
\end{tabular}
\end{center}

\subsection{Why This Works}

The mechanism is the Propagation Theorem applied at the observation level.  CEGIS operates on \texttt{PredObs} = \texttt{Carrier $\times$ Bool}.  It has no notion of \emph{which action pair} generated an observation.  When \texttt{prefer(A,B)} and \texttt{prefer(C,D)} are governed by the same \texttt{PredProg}, their observations are \emph{identical}:
\[
\texttt{(s, eval (feat f₁) s)} \quad\text{regardless of whether we asked about A-vs-B or C-vs-D.}
\]
Any \texttt{PredProg} evaluation is automatically feature-respecting (a direct corollary of the Propagation Theorem):

\begin{spec}
eval-respects : ∀ (p : PredProg) → FeatureRespects (λ s → eval p s)
eval-respects p c₁ c₂ afa =
  propagation p c₁ c₂ (all-agree→feat-equiv c₁ c₂ afa p)
\end{spec}

\noindent This guarantees that the cross-pair effect composes with dissolution: once a pair observation propagates across an equivalence class, it propagates across \emph{all} correlated ranking pairs simultaneously.

\begin{remark}[Scaling]
With $k$ actions, there are $\binom{k}{2}$ ranking pairs.  If all pairs are governed by the same feature structure (as when the total ordering is feature-determined), observing any single pair constrains all $\binom{k}{2}$ pairs.  Combined with dissolution ($m$ states per class $\to$ 1 observation), the savings grow as $\binom{k}{2} \times m$.  For $k = 4$ and $m = 2$: a factor of $18\times$.  For $k = 10$ and $m = 100$: a factor of $4{,}500\times$.
\end{remark}

\begin{remark}[Partial Correlation]
In general, not all ranking pairs will share the same predicate.  If \texttt{prefer(A,B)} = \texttt{feat~f₁} but \texttt{prefer(A,C)} = \texttt{feat~f₂}, then C-vs-D observations (which reveal $f_1$) give A-vs-B for free but not A-vs-C.  The propagation effect is strongest when the ranking structure has high \emph{feature correlation} across pairs---a property that can be checked at synthesis time by inspecting the surviving programs.
\end{remark}

\subsection{Verified Learning Curves}

To visualize the cross-pair effect, we instrument a richer scenario: 4 actions, 3 features ($f_1$, $f_2$, $f_3$), and 8 states (one per equivalence class).  The ranking structure has \emph{partial} correlation: $\mathsf{prefer}(A,B) = \mathsf{prefer}(C,D) = \mathsf{feat}~f_1$, $\mathsf{prefer}(A,C) = \mathsf{prefer}(B,D) = \mathsf{feat}~f_2$, and $\mathsf{prefer}(B,C) = \mathsf{feat}~f_3$.

We feed C-vs-D observations to CEGIS and track \emph{both} the version-space size and the fraction of states where the A$>$B ranking is correctly determined.  Every entry in Table~\ref{tab:curves} is verified by \texttt{refl} in Agda (module \texttt{PropagationCurves}).

\begin{table}[H]
\centering
\caption{Version-space evolution and A$>$B ranking accuracy under C-vs-D observations.  All values verified by \texttt{refl}.}
\label{tab:curves}
\small
\begin{tabular}{@{}clccl@{}}
\textbf{Step} & \textbf{Observation} & $\boldsymbol{|VS|}$ & \textbf{Accuracy} & \textbf{Note} \\
\hline
0 & --- & 5 & 0/8 & initial \\
1 & $C > D$ at $s_1$ {\scriptsize(T,T,T)} & 4 & 1/8 & \texttt{falsep} eliminated \\
2 & $C > D$ at $s_2$ {\scriptsize(T,T,F)} & 3 & 2/8 & \texttt{feat~f₃} eliminated \\
3 & $C > D$ at $s_3$ {\scriptsize(T,F,T)} & 2 & 4/8 & \texttt{feat~f₂} eliminated \\
4 & $C > D$ at $s_4$ {\scriptsize(T,F,F)} & 2 & 4/8 & non-informative \\
5 & $D > C$ at $s_5$ {\scriptsize(F,T,T)} & 1 & 8/8 & \texttt{truep} eliminated \\
6 & $D > C$ at $s_6$ {\scriptsize(F,T,F)} & 1 & 8/8 & redundant \\
7 & $D > C$ at $s_7$ {\scriptsize(F,F,T)} & 1 & 8/8 & redundant \\
8 & $D > C$ at $s_8$ {\scriptsize(F,F,F)} & 1 & 8/8 & redundant \\
\end{tabular}
\end{table}

\noindent The A-vs-B observation sequence produces \emph{identical} \texttt{PredObs} values, verified in Agda:

\begin{spec}
cross-pair-identity : cd-obs ≡ ab-obs
cross-pair-identity = refl
\end{spec}

\noindent Figure~\ref{fig:cross-pair-curves} visualizes the accuracy curve.

\begin{figure}[H]
\centering
\begin{tikzpicture}
\begin{axis}[
    width=0.88\textwidth,
    height=7cm,
    xlabel={Observations fed to CEGIS},
    ylabel={A$>$B ranking accuracy (\% of states correct)},
    xmin=-0.3, xmax=8.5,
    ymin=-5, ymax=115,
    xtick={0,1,...,8},
    ytick={0,12.5,25,50,100},
    yticklabels={0\%,12.5\%,25\%,50\%,100\%},
    grid=major,
    legend style={at={(0.02,0.98)}, anchor=north west, font=\small,
                  draw=nord3!50, fill=white, fill opacity=0.9},
    axis line style={nord3},
    tick style={nord3},
    grid style={nord3!20},
]
% A>B direct observations (blue, filled circles)
\addplot[nord10, very thick, mark=*, mark size=3pt] coordinates {
    (0,0) (1,12.5) (2,25) (3,50) (4,50) (5,100) (6,100) (7,100) (8,100)
};
\addlegendentry{From A$>$B observations (direct)}

% C>D cross-pair observations (red, triangles — IDENTICAL curve)
\addplot[nord11, thick, mark=triangle*, mark size=3.5pt, dashed,
         mark options={solid}] coordinates {
    (0,0) (1,12.5) (2,25) (3,50) (4,50) (5,100) (6,100) (7,100) (8,100)
};
\addlegendentry{From C$>$D observations (cross-pair)}

% Annotations
\node[above right, font=\scriptsize, text=nord3] at (axis cs:3.6,53)
    {non-informative};
\draw[->, nord3, thin] (axis cs:3.9,52) -- (axis cs:4,50.5);

\node[above, font=\scriptsize\bfseries, text=nord14] at (axis cs:5,105)
    {converged};
\draw[->, nord14, thin] (axis cs:5,104) -- (axis cs:5,101);

\draw[decorate, decoration={brace, amplitude=4pt, mirror}, nord3, thick]
    (axis cs:5.8,0) -- (axis cs:8.2,0)
    node[midway, below=5pt, font=\scriptsize, text=nord3] {all redundant};

\end{axis}
\end{tikzpicture}
\caption{A$>$B ranking accuracy as C-vs-D (sub-optimal) observations accumulate.  The blue circles (direct A$>$B observations) and red triangles (C$>$D cross-pair observations) produce \textbf{identical curves}---they overlap exactly.  CEGIS receives the same \texttt{PredObs} values regardless of which pair was observed.  Step~4 is non-informative (both surviving programs agree at that state).  After step~5 the version space has converged; steps~6--8 are redundant.  Every data point is verified by \texttt{refl} in Agda.}
\label{fig:cross-pair-curves}
\end{figure}

\noindent The plot tells the story directly: an agent that only ever observes the sub-optimal pair (C~vs~D) learns the top-tier ranking (A~$>$~B) with \emph{exactly} the same speed and accuracy as one that directly observes A~vs~B.  The curves don't merely correlate---they are definitionally equal.

\subsection{Rollout Performance at Scale}
\label{sec:rollout}

The verified curves above use 8~states.
To demonstrate the effect at a scale closer to a realistic decision-making problem, we instrument a larger environment:
\textbf{30~states}, \textbf{4~actions} (A--D), and \textbf{4~features} ($f_1{=}s{<}20$, $f_2{=}s\bmod2{=}0$, $f_3{=}s{<}10$, $f_4{=}s\bmod3{=}0$).

Arm rewards follow a contextual structure:
$r(s,\text{A}){=}5\,f_1{+}5\,f_2$, $r(s,\text{B}){=}5\,f_1{+}5\,(1{-}f_2)$,
$r(s,\text{C}){=}5\,(1{-}f_1){+}5\,f_2$, $r(s,\text{D}){=}5\,(1{-}f_1){+}5\,(1{-}f_2)$.
This gives four state--feature classes, each with a unique optimal arm, and an optimal total reward of~500 over a 50-step trajectory.

The \textbf{ranking structure} is:
A${>}$B and C${>}$D are governed by $\mathsf{feat}\;f_2$~(cross-pair group);
A${>}$C and B${>}$D are governed by $\mathsf{feat}\;f_1$~(cross-pair group);
A${>}$D and B${>}$C require depth-1 programs (unresolvable at depth~0).
The version space at depth~1 contains \textbf{22~programs} (literals, negations, conjunctions, disjunctions).

\subsection{Demo 13: FrozenLake \texorpdfstring{$4 \times 4$}{4x4} (Slippery) --- Stochastic Synthesis}
\label{sec:e2e-frozenlake-slippery}

Demos~7--12 operated in the \emph{deterministic} setting: actions have a single, predictable outcome.
We now cross the deterministic/stochastic boundary by implementing the \textbf{slippery FrozenLake}, a stochastic variant of the standard OpenAI Gym benchmark (\texttt{is\_slippery=True}).
The grid and map topology are identical to Demo~7 (same 16~states, same 4~holes and 1~goal), but each chosen action has only $\tfrac{1}{3}$ probability of executing in the intended direction, with $\tfrac{1}{3}$ each for the two perpendicular directions (e.g., choosing Right yields \{Right, Down, Up\}, each with weight~1 out of~3).

This demo achieves five firsts:
\begin{enumerate}
  \item The first \emph{stochastic RL benchmark} in the synthesis pipeline---\texttt{StochasticFiniteMDP} is instantiated with a real Gym-standard environment.
  \item A \textbf{memoized stochastic Finder} (\texttt{MemoizedStochastic}) builds a DP value table layer-by-layer, achieving $O(\text{depth} \times |S| \times |A| \times |\text{outcomes}|)$ total cost---polynomial in all parameters.
  \item The memoized Finder at \textbf{depth~10} produces a \emph{safety-aware} policy where all 4~actions appear with genuine expected-value justifications.
  \item \textbf{Decision-tree CEGIS} synthesizes \texttt{PredProg} policies for the stochastic environment from row/col features.
  \item The synthesized policy is \textbf{verified} to match the Finder at all 11~non-terminal states.
\end{enumerate}
The pipeline scales to the 8×8 grid (Demo~14, \S\ref{sec:e2e-frozenlake-8x8-slippery}).

\paragraph{Slippery step function.}
The step function uses the \texttt{Dist} type from \texttt{Probability.Finite}: for non-terminal states, it returns three equally-weighted outcomes (intended direction plus two perpendicular directions).
A key design choice: terminal states (holes and goal) also return \emph{three copies} of the self-loop rather than a single \texttt{pure} outcome.
This ensures a uniform branching factor of~3 at every state, which is critical for the memoized DP Finder: without uniform branching, the unnormalized expected values grow at rate $3^k$ for non-terminal states but only $k$ for terminal states, making comparisons across branches with mixed terminal/non-terminal successors incorrect.

\paragraph{Memoized stochastic Finder.}
The recursive \texttt{best-expected-trace} computes expected values by expanding the distribution tree, yielding exponential cost in the horizon.
We add a \texttt{MemoizedStochastic} module to the \texttt{StochasticFiniteMDP} EC that builds a value table bottom-up:
\[
  V_0(s) = 0, \quad V_{k+1}(s) = \max_a \bigl[\mathbb{E}[r \mid s,a] + \textstyle\sum_i w_i \cdot V_k(s'_i)\bigr]
\]
Each level reuses the previous table, so the total cost is $O(\text{depth} \times |S| \times |A| \times |\text{outcomes}|)$.
At depth~10 with 16~states, 4~actions, and 3~outcomes, the memoized Finder computes the entire policy table in under 20~seconds---compared to the recursive Finder, which was already strained at horizon~6.

\paragraph{Policy structure.}
At depth~10, the memoized Finder produces a genuinely \emph{safety-aware} policy with all four actions represented:
\begin{itemize}
  \item \textbf{D} (progress) at states \{0, 9, 14\}: these states can safely advance toward the goal column.
  \item \textbf{L} (safe retreat) at states \{4, 6, 10\}: row~1--2 states that avoid nearby holes by retreating left.
  \item \textbf{U} (hole avoidance) at states \{1, 3, 8\}: states adjacent to holes that stay in the safer upper rows.
  \item \textbf{R} (lateral progress) at states \{2, 13\}: states that advance rightward toward the goal column.
\end{itemize}
This is a qualitative improvement over the horizon-6 policy, where distant states were dominated by tie-breaking (L won by action ordering) rather than genuine expected-value differences.
The depth-10 policy resolves ties through accumulated probability mass---e.g., state~1 now chooses U (staying away from hole~5) instead of D (the tie-broken default at horizon~6).

\paragraph{Decision-tree synthesis.}
With 8~positional features (row-is~0--3, col-is~0--3), we cascade \texttt{synth-decision-tree} over four action classes: D~first (3~states positive, 8~negative), then L~among the remainder (3~positive, 5~negative), then U~(3~positive, 2~negative), with R~as default.
Each tree uses the depth-0 predicate pool (18~candidates) and fuel~8.

\paragraph{Verification.}
The synthesized policy is checked against the precomputed Finder map at all 11~non-terminal states:
\begin{center}
\texttt{test-auto-all-match : check-all finder-map $\equiv$ true}\\
\texttt{test-auto-all-match = refl}
\end{center}
Compile time: ${\sim}50$~seconds (memoized Finder + decision-tree synthesis).

\begin{remark}[Stochastic Synthesis is EC-Agnostic]
The synthesis infrastructure (\texttt{PredicateDSL}, \texttt{CEGIS}, \texttt{synth-decision-tree}) is completely independent of the transition model.
The same \texttt{FDMDPSynthesis} module that synthesizes deterministic FrozenLake policies also synthesizes stochastic ones---only the \emph{oracle} differs.
In the deterministic case, the oracle is a DP Finder; in the stochastic case, it is a memoized expected-value Finder.
The predicate synthesis never sees the step function; it only sees state--label pairs.
This EC-agnosticism is a direct consequence of the \texttt{PredicateDSL} operating on carriers, not transitions.
\end{remark}

\begin{remark}[Uniform Branching Invariant]
The memoized DP Finder computes unnormalized expected values: $V_k(s) = \sum_i w_i \cdot V_{k-1}(s'_i)$ without dividing by total weight.
For environments where all states have the same number of equally-weighted outcomes (uniform branching factor~$b$), unnormalized values are proportional to true expectations with a common factor $b^k$, so action rankings are preserved.
When terminal states self-loop via \texttt{pure} (branching factor~1) while non-terminal states branch (factor~3), the unnormalized values grow at different rates ($1^k$ vs $3^k$), corrupting comparisons.
The uniform branching design---where terminal states emit 3~copies of the self-loop---eliminates this issue structurally, without requiring rational arithmetic or explicit normalization.
\end{remark}

\subsection{Demo 14: FrozenLake \texorpdfstring{$8 \times 8$}{8x8} (Slippery) --- Scaling the Stochastic Pipeline}
\label{sec:e2e-frozenlake-8x8-slippery}

We scale the slippery FrozenLake to the 8×8 grid (64~states, 53~non-terminal, 10~holes, 1~goal), using the same memoized stochastic Finder at depth~10.
The map topology matches the standard OpenAI Gym 8×8 layout; movement uses arithmetic (\texttt{s/8}, \texttt{s\%8}) for scalability.
Uniform branching factor~3 is preserved at all states.

The memoized Finder computes the policy in ${\sim}27$~seconds, demonstrating that the polynomial DP scales to 4× the state count of the 4×4 variant.
Key policy probes: state~0 chooses L (safety at the start); states~56--57 choose D (progress along the bottom row); states~60--62 near the goal reflect the depth-10 expected-value trade-offs.
This is the first \emph{stochastic} RL benchmark at 64~states with verified policy computation in the synthesis pipeline.

\paragraph{Decision-tree synthesis.}
With 53~non-terminal states and four actions (D, L, U, R), the policy distribution is scattered: no single row or column feature isolates all states preferring a given action.
\texttt{synth-greedy-or} fails here because it builds flat disjunctions of atoms---each atom must have zero false positives---and with only \texttt{row-is} and \texttt{col-is}, many states sharing a row or column have different optimal actions.
We use \texttt{synth-decision-tree} instead: it greedily splits on the feature that minimizes Gini impurity, then recurses on each branch, building conjunctions implicitly (e.g., \texttt{row-is~7~$\wedge$~col-is~6} identifies state~62).
With 16~features (\texttt{row-is}~0--7, \texttt{col-is}~0--7) and depth-0 predicate pool, we cascade \texttt{synth-decision-tree} over four action classes: D~first, then L, then U, with R~as default.
Each stage uses fuel~8.

\paragraph{Verification.}
The synthesized policy is checked against the Finder map at all 53~non-terminal states:
\begin{center}
\texttt{test-auto-all-match : check-all finder-map $\equiv$ true}\\
\texttt{test-auto-all-match = refl}
\end{center}
Compile time: ${\sim}2.9$~hours (memoized Finder + decision-tree synthesis over 53~states).

\begin{center}
\begin{tabular}{@{}lcc@{}}
& \textbf{Demo 13 (4×4)} & \textbf{Demo 14 (8×8)} \\
\hline
States (non-terminal) & 16 (11) & 64 (53) \\
Features & 8 (row/col 0--3) & 16 (row/col 0--7) \\
Cascade & D$\to$L$\to$U$\to$R & D$\to$L$\to$U$\to$R \\
CEGIS & \texttt{synth-decision-tree} & \texttt{synth-decision-tree} \\
Finder time & ${\sim}$20 s & ${\sim}$27 s \\
Compile time & ${\sim}$50 s & ${\sim}$2.9 h \\
\end{tabular}
\end{center}

